% ============================================================
%  Chapter 7: Evaluation
%  Included from report.tex via % ============================================================
%  Chapter 7: Evaluation
%  Included from report.tex via % ============================================================
%  Chapter 7: Evaluation
%  Included from report.tex via % ============================================================
%  Chapter 7: Evaluation
%  Included from report.tex via \include{Chapters/Evaluation}.
%  Do NOT add \documentclass or \begin{document} here.
%  Follows Chapter 6 (Implementation); precedes Chapter 8
%  (Professional Issues / Conclusion).
%  Scope: tests every research question and pre-registered
%  hypothesis fixed in Chapter 4 against the run evidence written
%  by the system built in Chapter 6, including a comparison with
%  the alternative designs rejected in Chapter 5.
% ============================================================

\chapter{Evaluation}
\label{chap:evaluation}

This chapter asks whether the system specified in
Chapter~\ref{chap:spec} and built in Chapter~\ref{chap:impl} produces
the results predicted by the directional hypotheses fixed before any
run. Procedure and dataset are stated first
(\S\S\ref{sec:eval-procedure}--\ref{sec:eval-dataset}); each later
section reports against a target chosen before the data was seen.
Results are organised by research question
(\S\S\ref{sec:eval-headline}--\ref{sec:eval-rq3}), then by failure
mode and pipeline-design choice (\S\ref{sec:eval-failures}), then by
counterfactual against the alternatives discarded in
Chapter~\ref{chap:design} (\S\ref{sec:eval-alternatives}).
\S\ref{sec:eval-validity} states what these numbers cannot say;
\S\ref{sec:eval-synthesis} hands the transferable findings to
Chapter~\ref{chap:conclusion}.

% ------------------------------------------------------------
\section{Procedure and what counts as evidence}
\label{sec:eval-procedure}

Four runs were executed in April 2026, one per cell of the
$2\times 2$ factorial in Figure~\ref{fig:factorial}: Sonnet/OCR,
Sonnet/Vision, GPT-4o/OCR and GPT-4o/Vision. Each run processed all
forty documents in \texttt{data/dataset\_manifest.csv} with status
\texttt{active}, using the prompt template archived alongside the run
directory. Pipeline configuration, prompt and manifest snapshot were
written verbatim into each run directory at start time, so any later
run is a diff inspection of those three files away from the four
reported here.

The metric set is the one fixed in
Chapter~\ref{chap:spec}~\S\ref{sec:spec-evaluation}: field-level
accuracy with both-null cells excluded from the denominator,
document-level accuracy (a document counts as correct only if all
twelve fields are correct), and overall accuracy aggregated across
fields and documents. The four-class null taxonomy of
Figure~\ref{fig:null-taxonomy} is applied at field-comparison time
and persisted, so every later analysis below reads from the same
classification. Falsification follows the pre-registration: H1 is
rejected if vision is below OCR under the same provider, H2 is
rejected if any single condition is strictly cheaper, faster and
simpler than its peers, and H3 is rejected if the within-provider
modality gap is the same sign and magnitude across providers.

Three properties of the evidence are worth fixing now. Every run
completed end-to-end: 0/40 documents failed in any condition, so no
result below is restricted to a successful subset. No JSON-parse
failure occurred across the $40\times 12\times 4 = 1{,}920$ field
comparisons, so the parse-hardening of
Chapter~\ref{chap:design}~\S\ref{sec:design-execution} held and the
content findings are not contaminated by formatting failures. The
corpus is small ($n=40$) and errors cluster at the document level,
so per-language slices are reported as descriptive ranges rather than
as point estimates with confidence intervals. Significance testing is
not attempted, in line with the scope decision in
\S\ref{sec:spec-scope}.

% ------------------------------------------------------------
\section{Dataset and ground truth}
\label{sec:eval-dataset}

\subsection{Provenance}
\label{sec:eval-dataset-provenance}

Utility bills are personal-data instruments under GDPR
Article~4(1)~\cite{eu2016gdpr}. Random sampling from a billing
provider's archive, the textbook way to build a balanced
multilingual corpus, is therefore not feasible at
undergraduate-project scale. The corpus had to be assembled by hand.

Two streams contribute. Stream~A is bills released publicly by the
issuers themselves: redacted exemplars on customer-help pages,
sample bills on tariff-board portals, and reproductions on
billing-explainer guides. Stream~A accounts for the majority of the
corpus and covers the British, American, Australian, Canadian,
French, German and Italian subsets. Stream~B is bills contributed
by personal contacts under written consent for use in this study;
Stream~B sits inside the British and Italian subsets, fills two gaps
where public specimens were thin, and was used only after every
contributor signed a one-page consent form covering scope and
storage. Every document in either stream had personally identifiable
account, customer and address detail redacted before storage in the
repository, in line with the provenance controls
Chapter~\ref{chap:lsesp} sets out.

The corpus aims for within-budget diversity rather than statistical
representativeness, and the implications for external validity
appear in \S\ref{sec:eval-validity}. The provenance constraints are
not a defect of the dataset; they are the reason the dataset has
the shape it has.

\subsection{Composition}
\label{sec:eval-dataset-composition}

Forty documents are in scope. Twenty are English, twelve are German,
seven are Italian and one is French. Twenty-seven are electricity
bills, six are gas, and seven are water. Thirty-three unique issuers
appear; concentration is tracked in \S\ref{sec:eval-validity}
because four BC~Hydro bills, two Affinity~Water bills, two
Con~Edison bills, two Vattenfall Europe Sales bills, two Green
Planet Energy bills, two Lupatotina bills, and two PG\&E bills carry
greater weight on any per-issuer slice than a single document would.
Page-count distribution: eleven documents have at most two pages,
and twenty-nine documents (72.5\% of the corpus) include at least
one page beyond two. Figure~\ref{fig:eval-dataset-composition} shows
the three breakdowns on one row; the page-count panel returns in
\S\ref{sec:eval-failures-pagecap} when the two-page baseline is
defended empirically.

\begin{figure}[ht]
    \centering
    \includegraphics[width=\linewidth]{Chapters/figures/fig-7-7-dataset-composition.png}
    \caption[Dataset composition.]{The forty-document corpus by
    language, by utility type and by page count. The page-count panel
    marks the two-page pipeline cap; documents to the right of the
    dashed line have at least one page that neither pipeline
    processes. The cap is defended in
    \S\ref{sec:eval-failures-pagecap}.}
    \label{fig:eval-dataset-composition}
\end{figure}

\subsection{Annotation and the meaning of ``correct''}
\label{sec:eval-dataset-annotation}

A single annotator (the author) produced the per-document JSON files
in \texttt{data/ground\_truth/} using the format fixed in
\S\ref{sec:spec-schema}. Dates were recorded in ISO~8601, numerics
in locale-preserved form, and \texttt{null} was reserved for fields
genuinely absent from the bill rather than for fields that were
unreadable. The four-class taxonomy of
Figure~\ref{fig:null-taxonomy} therefore maps cleanly to predictions
without further annotation steps. Where a field admitted multiple
defensible answers the annotation records one of them, and the
choice is itself an annotation decision rather than a fact about the
bill; the specific cases are listed in \S\ref{sec:eval-validity} and
their effect on the headline numbers is bounded there.
Single-annotator coverage is acknowledged in
\S\ref{sec:eval-validity} as well, because inter-annotator agreement
cannot be computed from a single labelling pass.

% ------------------------------------------------------------
\section{Headline results}
\label{sec:eval-headline}

\begin{table}[ht]
    \centering
    \caption{Field-level overall accuracy and document-level accuracy
    across the four conditions. Both-null cells are excluded from the
    denominator (\S\ref{sec:spec-evaluation}). Source:
    \texttt{results/reports/four\_way\_20260429/comparison.json}.}
    \label{tab:eval-headline}
    \begin{tabular}{lcc}
        \hline
        Condition & Field-level overall & Document-level \\
        \hline
        Sonnet 4.5 / OCR-mediated     & $89.65\%$ & $45.0\%$ \\
        Sonnet 4.5 / Vision-mediated  & $87.55\%$ & $37.5\%$ \\
        GPT-4o / OCR-mediated         & $86.35\%$ & $30.0\%$ \\
        GPT-4o / Vision-mediated      & $85.33\%$ & $22.5\%$ \\
        \hline
    \end{tabular}
\end{table}

Table~\ref{tab:eval-headline} reports the four overall numbers and
the four document-level numbers. Figure~\ref{fig:eval-overall} is
the bar-chart view; Figure~\ref{fig:eval-heatmap} is the per-field
heatmap. At field level the four conditions cluster within a
4.32-percentage-point band; at document level the band widens to
22.5 points because a document is correct only when all twelve
fields are correct, and single-field misses across the schema
accumulate. The document-level signal is the more
production-relevant of the two
(Chapter~\ref{chap:spec}~\S\ref{sec:spec-evaluation}): a bill with
one wrong field cannot be ingested without manual review. None of
the four conditions clears 50\% document-level accuracy, which is
itself a finding about the gap between aggregate accuracy and
end-to-end usefulness on this document class at this prompt and
this schema.

\begin{figure}[ht]
    \centering
    \includegraphics[width=0.85\linewidth]{Chapters/figures/fig-7-1-overall-accuracy.png}
    \caption[Overall accuracy, four conditions.]{Field-level overall
    accuracy across the four conditions. Sonnet/OCR is the highest
    cell; GPT-4o/Vision is the lowest. The band is narrow at field
    level and wide at document level (Table~\ref{tab:eval-headline}).}
    \label{fig:eval-overall}
\end{figure}

\begin{figure}[ht]
    \centering
    \includegraphics[width=\linewidth]{Chapters/figures/fig-7-2-field-accuracy-heatmap.png}
    \caption[Per-field accuracy heatmap.]{Per-field accuracy across
    twelve fields and four conditions. \texttt{utility\_type},
    \texttt{currency} and \texttt{consumption\_unit} are uniformly
    high; \texttt{bill\_number}, \texttt{account\_number} and
    \texttt{provider\_name} are where the four conditions disagree
    the most. Source: \texttt{accuracy\_matrix} in
    \texttt{comparison.json}.}
    \label{fig:eval-heatmap}
\end{figure}

% ------------------------------------------------------------
\section{RQ1 --- does input modality matter?}
\label{sec:eval-rq1}

\subsection{Aggregate verdict on H1}
\label{sec:eval-rq1-aggregate}

H1 (vision $\geq$ OCR within provider) is rejected on this corpus.
For Sonnet, OCR-mediated extraction beats vision-mediated extraction
by $+2.10\,\mathrm{pp}$ at field level (89.65\% vs.\ 87.55\%); for
GPT-4o the gap is $+1.02\,\mathrm{pp}$ (86.35\% vs.\ 85.33\%). Both
gaps point the wrong way for H1. The strict reading is that H1 is
falsified for both providers. The cautious reading is that for
GPT-4o the gap is within the variation expected from $n=40$ and the
finding is then that the modality effect is much smaller for GPT-4o
than the literature
predicted~\cite{shen2026ocrornot}. The prior that
vision-mediated extraction would dominate OCR-mediated extraction
does not survive contact with this corpus.

\subsection{Where the gap actually lives}
\label{sec:eval-rq1-stratified}

The aggregate gap conceals slice-level differences large enough to
matter for system design. Figure~\ref{fig:eval-by-lang} stratifies
overall accuracy by language. The
``OCR\,$>$\,Vision'' finding is dragged by English: for Sonnet,
$\Delta_{\mathrm{en}} = +3.42\,\mathrm{pp}$; for GPT-4o,
$\Delta_{\mathrm{en}} = +2.15\,\mathrm{pp}$. On Italian
($n=7$) the GPT-4o ranking reverses to
$-2.72\,\mathrm{pp}$ (vision wins), and on French ($n=1$) Sonnet/Vision
is exactly correct on every field.
Figure~\ref{fig:eval-by-utility} is the same picture by utility
type: the smallest slice (water, $n=7$) shows the largest reversal
under Sonnet (vision narrowly behind on water but ahead on gas),
again from inside an aggregate that says vision is behind.

\begin{figure}[ht]
    \centering
    \includegraphics[width=0.95\linewidth]{Chapters/figures/fig-7-3-accuracy-by-language.png}
    \caption[Accuracy by language.]{Field-level accuracy by language
    for the four conditions. The aggregate ``OCR\,$>$\,Vision''
    finding is dragged by English; on Italian the GPT-4o ranking
    reverses, and on French ($n=1$) Sonnet/Vision is perfect.}
    \label{fig:eval-by-lang}
\end{figure}

\begin{figure}[ht]
    \centering
    \includegraphics[width=0.95\linewidth]{Chapters/figures/fig-7-4-accuracy-by-utility-type.png}
    \caption[Accuracy by utility type.]{Field-level accuracy by
    utility type. Water bills are the smallest slice ($n=7$) but
    show non-trivial modality reversals; the largest gas-side gap is
    GPT-4o, where OCR beats vision by $7.0\,\mathrm{pp}$.}
    \label{fig:eval-by-utility}
\end{figure}

\subsection{Field-level offsets that point in opposite directions}
\label{sec:eval-rq1-flips}

The aggregate gap is composed of per-field offsets that pull both
ways inside the same documents. Figure~\ref{fig:eval-flips} shows
the count of documents where one modality is correct on a given
field and the other is wrong, broken out per provider. For Sonnet,
\texttt{provider\_name} flips OCR-only-correct on six documents and
Vision-only-correct on two, a six-to-two asymmetry that contributes
the largest single chunk of the modality gap; on
\texttt{total\_amount\_due} the asymmetry inverts (Vision-only-correct
on four documents, OCR-only-correct on one). For GPT-4o,
\texttt{account\_number} flips Vision-only-correct on eleven
documents against OCR-only-correct on one, an asymmetry that pulls
GPT-4o/Vision \emph{ahead} on that field even though it is behind in
aggregate; \texttt{provider\_name} pulls in the opposite direction
(OCR-only-correct on nine documents against Vision-only-correct on
two).

\begin{figure}[ht]
    \centering
    \includegraphics[width=0.95\linewidth]{Chapters/figures/fig-7-11-modality-flip-counts.png}
    \caption[Modality-flip counts per field and provider.]{Per
    provider, per field, the count of documents where exactly one
    modality is correct. The two columns ``OCR-only correct'' and
    ``Vision-only correct'' sum to disagreement volume on that field
    for that provider, regardless of whether the modality wins or
    loses on that field overall.}
    \label{fig:eval-flips}
\end{figure}

The flip counts are the empirical case against any single-modality
system on this corpus. Twelve documents in the corpus carry at least
one such flip per provider, and on several documents the flips
oppose: a perfect modality router would beat the better of the two
single-modality systems by several percentage points
(\S\ref{sec:eval-alternatives}).

\subsection{Mechanisms behind the gap}
\label{sec:eval-rq1-mechanisms}

Four mechanisms account for most of the failures that drive the
aggregate gap. Each is identified by inspecting the wrong
predictions on representative documents in
\texttt{results/runs/}; named documents are given so the
attribution is auditable.

\textbf{(M1) Brand wordmark vs.\ legal name on
\texttt{provider\_name}.} Sonnet/Vision drops from 0.85 to 0.75 on
this field; GPT-4o drops from 0.825 to 0.65. Vision systematically
prefers the prominent brand wordmark on the bill header
(``ENGIE'', ``IREN'', ``Xcel Energy'') to the legal-entity surface
form recorded as ground truth. The OCR-mediated path sees the legal
suffix in the small print at the bottom of the bill; the
vision-mediated path anchors to whichever string is the largest
piece of type on the page. This is a header-salience effect, not a
legal-suffix-stripping effect: the ground truth records one of two
defensible surface forms and the modality systematically picks
different ones.

\textbf{(M2) Whitespace coercion on \texttt{account\_number}.} GPT-4o
in OCR mode scores 0.486 on this field while Sonnet in OCR mode
scores 0.750 and GPT-4o in vision mode scores 0.750. Inspection of
the GPT-4o/OCR misses shows GPT-4o stripping internal spaces from
account numbers it received with spaces in the OCR text, then
mismatching the spaced ground truth. Sonnet preserves the
whitespace; vision sees the layout directly and reproduces it. This
is a within-provider modality reversal: GPT-4o/Vision beats
GPT-4o/OCR on this field by 26.4~percentage points.

\textbf{(M3) Locale-aware separator handling on
\texttt{consumption\_amount}.} Both providers drop on this field
under vision: Sonnet from 0.923 to 0.846, GPT-4o from 0.923 to
0.821. The recurring failure pattern is European thousands
separators (``\texttt{16.250}'', ``\texttt{16\,250}'') being read by
the vision pipeline as a decimal (``\texttt{16.25}''). The
OCR-mediated path passes the literal string through the EU/US-aware
float parser of \S\ref{sec:design-normalisation}, which resolves
the locale at parse time using the per-document language flag.
Vision pipelines do their own number reading inside the model and
do not see that flag. The mechanism predicts that vision will
under-perform on European-locale numerics, which is the slice
pattern Figure~\ref{fig:eval-by-lang} shows.

\textbf{(M4) Field disambiguation on documents with multiple
candidates.} On bills that present several plausible totals
(``Current charges'', ``Outstanding balance'', ``Amount due''),
several plausible reference numbers (``Bill number'', ``Account
number'', ``Statement reference''), or several dates (``Bill date'',
``Issued'', ``Statement date''), both modalities sometimes pick the
wrong one. The clearest case is
\texttt{gb\_water\_affinitywater\_001}, where the ground-truth
\texttt{total\_amount\_due} is \pounds46.00 (current charges) but the
bill prominently displays \pounds235.14 (outstanding balance). The
disambiguation is field-construct ambiguity rather than a modality
problem: it appears in both modalities and is annotated as a
construct threat in \S\ref{sec:eval-validity}.

% ------------------------------------------------------------
\section{RQ2 --- the operational trade-off}
\label{sec:eval-rq2}

\subsection{Cost}
\label{sec:eval-rq2-cost}

H2 in its as-written form (vision more expensive) is supported.
Vision costs more per document for both providers: Sonnet/Vision
\$0.0153 against Sonnet/OCR \$0.0119 ($+28.8\%$), and GPT-4o/Vision
\$0.00938 against GPT-4o/OCR \$0.00828 ($+13.2\%$). Both pipelines
process the same first two pages of each bill (the design choice
defended in \S\ref{sec:eval-failures-pagecap}), so cost differences
trace directly to input-token volume: image tokens for the two
rasterised pages cost more than the text tokens for the
Tesseract-extracted text of the same two pages, on both providers'
2026 price lists. Figure~\ref{fig:eval-cost} is the per-document
view summed across the corpus.

\begin{figure}[ht]
    \centering
    \includegraphics[width=0.85\linewidth]{Chapters/figures/fig-7-5-cost-comparison.png}
    \caption[Per-document cost.]{Per-document estimated cost across
    the four conditions, summed over forty documents. Vision is
    more expensive than OCR for both providers; the gap is larger
    for Sonnet than for GPT-4o.}
    \label{fig:eval-cost}
\end{figure}

\subsection{Latency}
\label{sec:eval-rq2-latency}

H2 in its as-written form (vision slower) is rejected. Vision is
faster end-to-end for both providers. Mean wall-clock per document
is 9.65~seconds for GPT-4o/Vision against 14.44~seconds for
GPT-4o/OCR ($-33\%$), and 12.16~seconds for Sonnet/Vision against
16.35~seconds for Sonnet/OCR ($-26\%$). Figure~\ref{fig:eval-timing}
splits the time. Tesseract accounts for $\sim$8.7~seconds of every
OCR-mediated run; LLM-call latency rises in vision mode (Sonnet/Vision
8.6~seconds vs.\ Sonnet/OCR 3.7~seconds; GPT-4o/Vision 6.2~seconds
vs.\ GPT-4o/OCR 2.1~seconds) because the model now does the page
reading internally, but removing Tesseract is the larger effect on
total time on this hardware.

\begin{figure}[ht]
    \centering
    \includegraphics[width=0.85\linewidth]{Chapters/figures/fig-7-6-timing-breakdown.png}
    \caption[Timing breakdown.]{Mean per-document timings split into
    OCR and LLM stages. Tesseract dominates the OCR-mediated
    pipeline on this hardware; the LLM-call cost rises in vision
    mode but does not close the gap.}
    \label{fig:eval-timing}
\end{figure}

The latency ranking is hardware-conditioned. Tesseract was run
single-process on commodity CPU and is embarrassingly parallel, so
the OCR-side latency would shrink substantially on a multi-core
deployment. The latency ranking above is therefore a single-laptop
ranking, not a production-scale verdict; the production-scale
prediction is that the LLM-call dominates wall-clock under either
modality once OCR is parallelised, and the cross-modality latency
gap closes.

\subsection{Operational complexity}
\label{sec:eval-rq2-complexity}

Vision is simpler to deploy than OCR, and the ranking is
unambiguous. The OCR-mediated pipeline depends on Poppler (system
binary), Tesseract (system binary plus per-language data files for
\texttt{eng}, \texttt{deu}, \texttt{fra} and \texttt{ita}) and the
EU/US-aware float parser of
\S\ref{sec:design-normalisation}. The vision-mediated pipeline
depends on the provider SDK alone. Adding a fifth language to the
OCR-mediated pipeline requires a new Tesseract data file on every
host; adding a fifth language to the vision-mediated pipeline
requires nothing.

The H2 verdict is therefore mixed: vision is more expensive but
faster and operationally simpler. The trade-off direction H2
predicted is right; the per-axis directions split along provider
lines and along the cost-vs-latency axis in ways the as-written
hypothesis did not anticipate.

% ------------------------------------------------------------
\section{RQ3 --- does the choice of LLM matter?}
\label{sec:eval-rq3}

\subsection{Modality-by-provider interaction}
\label{sec:eval-rq3-interaction}

H3 (the modality gap differs by provider) is supported. The
within-provider modality gap is $+2.10\,\mathrm{pp}$ for Sonnet
and $+1.02\,\mathrm{pp}$ for GPT-4o, a gap-of-gaps of
$1.08\,\mathrm{pp}$. The interaction has a mechanistic reading:
Sonnet has more material to lose in vision mode because it is the
better OCR-mode reader on this corpus, and Sonnet's vision-mode
disadvantage concentrates on the M1 (header anchoring) and M3
(locale formatting) mechanisms above. GPT-4o is less prone to
header anchoring on the M1 cases inspected, partly because its
visible behaviour is to omit a field rather than substitute a
salient nearby string when the answer is uncertain
(\S\ref{sec:eval-rq3-orthogonal}).

\subsection{Provider effects orthogonal to modality}
\label{sec:eval-rq3-orthogonal}

Two provider effects appear in both modalities and are therefore
attributable to the model rather than to the input mode.

The first is currency contextual inference. The four BC~Hydro bills
have \texttt{currency} ground truth \texttt{CAD}, but the string
``CAD'' does not appear in the OCR text or on the page image: the
bills display the bare ``\$'' sign. Sonnet returns \texttt{CAD} on
all four under both OCR and vision; GPT-4o returns \texttt{USD} on
all four under both modalities. The signal Sonnet uses (``BC
Hydro'', ``Vancouver'', Canadian postal code) is available to both
providers; GPT-4o reads it and does not act on it. This single
mechanism accounts for the entire 0.90 vs.\ 1.00 currency gap in
Table~\ref{tab:eval-headline} between GPT-4o and Sonnet.

The second is the hallucination-versus-omission profile. Across
both modalities, Sonnet emits 20 hallucinations (8 OCR, 12 vision)
and 20 omissions (9 OCR, 11 vision); GPT-4o emits 5 hallucinations
(1 OCR, 4 vision) and 32 omissions (17 OCR, 15 vision). Sonnet
hallucinates roughly four times as often as GPT-4o; GPT-4o omits
where Sonnet guesses. Figure~\ref{fig:eval-pareto} plots the four
conditions on the two axes and shows that the providers occupy
distinct points, not different points on a single curve. Downstream
systems that punish wrong values should prefer GPT-4o; ones that
punish missing values should prefer Sonnet.

\begin{figure}[ht]
    \centering
    \includegraphics[width=0.78\linewidth]{Chapters/figures/fig-7-9-hallucination-omission-pareto.png}
    \caption[Hallucination-versus-omission Pareto.]{Per-condition
    counts of hallucinations and omissions across the corpus.
    Sonnet sits in the high-hallucination low-omission region under
    both modalities; GPT-4o sits in the low-hallucination
    higher-omission region. The two providers occupy distinct
    operating points, not different points on a single curve.}
    \label{fig:eval-pareto}
\end{figure}

\subsection{Cross-provider agreement and ensemble ceiling}
\label{sec:eval-rq3-agreement}

The two providers agree on the same answer for $\sim$80.4\% of OCR-mode
field comparisons and $\sim$78.1\% of vision-mode comparisons, where
``agree'' means both predictions normalise to the same value
regardless of correctness. Disagreements concentrate on the
mechanisms above: BC~Hydro \texttt{currency}, German credit-bill
totals, the Affinity Water \texttt{total\_amount\_due}
disambiguation, and \texttt{provider\_name} on bills with strong
brand wordmarks. A field-level vote ensemble of Sonnet and GPT-4o
in OCR mode would clear $\sim$91.4\% accuracy on this corpus
(against the better of the two cells alone, 89.65\% for Sonnet/OCR);
in vision mode the corresponding upper bound is $\sim$89.6\% against
87.55\% for Sonnet/Vision. The provider-level personality difference
H3 predicts is therefore exploitable, but the ceiling is not large
enough to outweigh the cost of running both providers per document.

% ------------------------------------------------------------
\section{Where failures actually come from}
\label{sec:eval-failures}

\subsection{The diagnosis heuristic and what it misses}
\label{sec:eval-failures-heuristic}

The diagnosis heuristic specified in
\S\ref{sec:design-evalframework-diagnosis} attributes a non-correct
OCR-mediated field to OCR if the ground-truth value is not a
case-insensitive substring of the stored OCR text, and to the LLM
otherwise. The heuristic is reproducible and inspectable, and it is
the right starting point. It also has known failure modes that
matter for the analysis below.

The first is contextual inference. On the BC~Hydro currency cases
above, the ground-truth string ``CAD'' is not in the OCR text, so
the strict heuristic attributes the GPT-4o failure to OCR. The
actual cause is that GPT-4o did not infer the currency from the
issuer name and address; Sonnet did infer it from exactly the same
OCR text. Attributing this to OCR is wrong in the sense that no OCR
change would have fixed it. The second is ground-truth phrasing: on
\texttt{us\_electricity\_publicservicecompany\_001} the GT
\texttt{provider\_name} is the legal name ``Public Service Company
of Colorado'', and the OCR text contains the brand wordmark ``Xcel
Energy'' rather than the legal name. The strict heuristic attributes
this to OCR; the actual cause is annotation choosing the legal name
where the brand wordmark would have been an equally defensible
answer. The third is locale-formatted numerics: on a bill that
displays ``\texttt{16.250}'' the GT records ``\texttt{16250}''
(thousands separator removed) and the heuristic again says OCR even
though the text was readable.

\subsection{Two-bucket vs.\ four-bucket re-attribution}
\label{sec:eval-failures-fourbucket}

Replacing two buckets with four addresses the failure modes above.
The four buckets are: \emph{Pure OCR} (Tesseract did not capture
the value), \emph{Pure LLM-extraction} (the LLM mis-read text that
was present), \emph{LLM-inference} (context for inference was
present and the model did not act on it; the BC~Hydro currency
cases sit here), and \emph{GT-ambiguity} (the annotation chose one
of multiple defensible answers; the PSCo and Affinity Water cases
sit here). The re-classification rule is: known multiple-answer
cases move to GT-ambiguity, currency and \texttt{utility\_type}
strict-OCR-failures move to LLM-inference, the rest stay where the
strict heuristic put them. Figure~\ref{fig:eval-attribution} shows
the result.

\begin{figure}[ht]
    \centering
    \includegraphics[width=0.95\linewidth]{Chapters/figures/fig-7-10-failure-attribution.png}
    \caption[Failure attribution: two-bucket vs.\ four-bucket.]{Failure
    attribution for the two OCR-mediated runs. Panel (a): the
    strict two-bucket heuristic of
    \S\ref{sec:design-evalframework-diagnosis}. Panel (b): four
    buckets after applying the re-classification rule of
    \S\ref{sec:eval-failures-fourbucket} to
    \texttt{derived/diagnosis\_audit.csv}.}
    \label{fig:eval-attribution}
\end{figure}

Of 47 Sonnet/OCR non-correct fields, 26 (55\%) are Pure OCR, 17
(36\%) are Pure LLM-extraction, 1 is LLM-inference and 3 are
GT-ambiguity. Of 61 GPT-4o/OCR non-correct fields, 31 (51\%) are
Pure OCR, 23 (38\%) are Pure LLM-extraction, 4 are LLM-inference
and 3 are GT-ambiguity. The headline finding survives in spirit: a
little over half of OCR-mode failures are addressable by improving
Tesseract; the remainder are addressable by improving the LLM, the
prompt or the annotation. The four-bucket reading sharpens that
finding by separating the cases where the OCR text was the problem
from the cases where the OCR text was fine and the model did not
act on what it had.

\subsection{Modality-specific patterns}
\label{sec:eval-failures-modality}

Two failure types appear under one input mode and not the other.
The OCR-mediated path hallucinates redacted spans into plausible
characters: on
\texttt{it\_electricity\_irenmercato\_001} the ground-truth
\texttt{bill\_number} is \texttt{502305934151}; the OCR text
contains the redacted form ``\texttt{50230XXXXXXX}'' and the LLM
returns a guessed completion in OCR mode while correctly returning
\texttt{null} (or the unredacted prefix) in vision mode where the
redaction box is visible as a layout element. The vision-mediated
path inverts the failure mode: on bills with prominent brand
wordmarks the model anchors to the wordmark on
\texttt{provider\_name} (M1), and on bills with European-locale
thousands separators the model collapses the separator (M3).

A third pattern affects every condition equally: sign handling on
credit and refund bills. On \texttt{de\_electricity\_naturstromhandel\_001}
the ground-truth \texttt{total\_amount\_due} is $-40.14$~EUR
because the bill is a refund (``Guthaben''). All four conditions
return $40.14$, dropping the sign. This is not a modality
strength on either side; it is a construct issue between the prompt
and the ground-truth convention, and it is recorded as a
limitation in \S\ref{sec:eval-validity}.

\subsection{The two-page baseline as a deliberate design choice}
\label{sec:eval-failures-pagecap}

Both pipelines process only the first two pages of each bill. The
cap is set explicitly in \texttt{src/pipeline.py} (lines 137 and
149: \texttt{ocr\_images = images[:2]} and
\texttt{vision\_images = images[:2]}). It is a deliberate baseline,
not a confound. On every bill in this corpus, every required field
that the bill carries appears within the first two pages: the
header (provider, utility type, bill number, dates), the totals
panel and the consumption box are universally on page one or two,
and later pages carry per-meter detail, terms and conditions,
multilingual translations or marketing inserts that the schema does
not target. The cap therefore costs almost nothing in accuracy
under this corpus while cutting input tokens by a factor of
$\sim$3.9 (the mean page count is 3.9 with the longest bill at
nine pages) and removing $\sim$5--7 seconds of Tesseract time on
the OCR side per long bill. Figure~\ref{fig:eval-page-cap}
visualises which documents are seen in full and which are
truncated.

\begin{figure}[ht]
    \centering
    \includegraphics[width=0.78\linewidth]{Chapters/figures/fig-7-8-page-cap-coverage.png}
    \caption[Two-page cap coverage.]{Per-document coverage under
    the two-page cap, applied to both pipelines. Dark segments are
    the pages either pipeline reads; light segments are the pages
    neither pipeline reads.}
    \label{fig:eval-page-cap}
\end{figure}

The cap is a property of the input distribution, not of the
algorithm. On a corpus that exercises later-page targets, changing
the slice from \texttt{[:2]} to a higher value, or to a function
that picks pages by a header-detection heuristic, would extend
coverage at proportional cost. The cost-vs-coverage curve for that
extension is the most direct empirical follow-up the corpus
suggests, and is left for Chapter~\ref{chap:conclusion}. The
relevant point here is that the cap is the reason cost and latency
are flat in document length on this corpus, and the reason the
modality-vs-modality comparison in \S\ref{sec:eval-rq1} is a
like-for-like comparison rather than a coverage-difference artefact.

% ------------------------------------------------------------
\section{Counterfactuals against the alternative designs}
\label{sec:eval-alternatives}

This section returns to the alternatives recorded in
Figure~\ref{fig:design-decisions} and tests them against the run
evidence rather than against intuition; the analysis closes the
A++ requirement that an evaluation chapter compare results expected
from theory for alternative designs with those obtained in
practice (Chapter~\ref{chap:spec}~\S\ref{sec:spec-evaluation}).

\subsection{Per-issuer templates}
\label{sec:eval-alternatives-templates}

The data ratifies the rejection of \S\ref{sec:design-decisions}.
Thirty-three unique issuers appear in the corpus and the field-level
error patterns are not positional: \texttt{provider\_name} fails on
brand-vs-legal selection (M1), \texttt{consumption\_amount} fails on
locale-aware separator handling (M3), and the
\texttt{total\_amount\_due} sign loss on credit bills is a semantic
question, not a coordinate-extraction question. None of those failure
modes is template-fixable; each one requires either an
issuer-specific rule or a flexible reading of the bill, and the
LLM-mediated approach gives the latter for free.

\subsection{Fine-tuned encoder--decoder (LayoutLMv3, Donut)}
\label{sec:eval-alternatives-finetune}

Forty documents is well below the floor at which Donut and
LayoutLMv3 are reported usable~\cite{kim2022donut,huang2022layoutlmv3};
provider concentration would inject training-distribution bias the
holdout cannot detect (BC~Hydro at four documents would be
disproportionately easy on any random split); and per-language
fine-tuning would be needed for the Italian and German subsets where
mechanisms M3 and the credit-bill sign issue occur. The LLM-mediated
approach pays a per-call cost in exchange for zero training data,
and at \$0.008--0.015 per document the per-call cost is below the
cost of one annotation hour at any reasonable rate.

\subsection{Provider ensembling}
\label{sec:eval-alternatives-ensemble}

The ensembling alternative is testable directly against the run
evidence (\S\ref{sec:eval-rq3-agreement}). A field-level vote of
Sonnet/OCR and GPT-4o/OCR clears $\sim$91.4\% on this corpus,
against 89.65\% and 86.35\% for the two cells alone. Most of the
gain is on the disagreement set (BC~Hydro \texttt{currency},
PSCo \texttt{provider\_name}, the German credit-bill totals where
Sonnet recovers the sign for one of the four cases). The ceiling
is real but moderate; pairing it with the cost ranking in
\S\ref{sec:eval-rq2-cost} (running two providers per document
roughly doubles the cost) puts the operating decision squarely in
the downstream-cost-tolerance domain.

\subsection{Single-modality system}
\label{sec:eval-alternatives-singlemod}

A vision-only or OCR-only system would have to cope with the
within-document modality flipping documented in
\S\ref{sec:eval-rq1-flips}. Twelve documents per provider have at
least one field where one modality is correct and the other wrong,
and several of those documents the directions oppose. A perfect
modality router would beat the better of the two single-modality
cells by roughly the same amount as a provider ensemble. The case
for keeping both pipelines is precisely this: not because either
wins outright, but because the conditions under which each is
preferable are field-, locale- and issuer-specific.

\subsection{Vision (or OCR) with all pages}
\label{sec:eval-alternatives-allpages}

The two-page cap is the design choice that makes both modalities
flat in document length on this corpus
(\S\ref{sec:eval-failures-pagecap}). Removing the cap on a corpus
where every required field already lives on the first two pages
would buy nothing in accuracy and cost roughly $3.9\times$ more in
input tokens (and proportionally more in OCR latency) on average.
Removing the cap on a different corpus where required fields appear
on later pages would buy accuracy at proportional cost, and the
cost-vs-coverage curve as a function of page count is the
empirical question Chapter~\ref{chap:conclusion} picks up. Either
way, the cap is the right choice for the corpus this study
evaluates.

% ------------------------------------------------------------
\section{What these numbers cannot say}
\label{sec:eval-validity}

\subsection{Internal validity}
\label{sec:eval-validity-internal}

Two internal threats matter. The diagnosis heuristic mis-attributes
the cases listed in \S\ref{sec:eval-failures-heuristic}; the
four-bucket re-classification of
\S\ref{sec:eval-failures-fourbucket} is the corrective. Sign
handling on credit-bill totals fails uniformly across all four
conditions because the prompt does not encode the
``\texttt{total\_amount\_due} is signed (negative on credit
bills)'' convention that the ground truth uses; this is a prompt
construct issue rather than a model failure, and it bounds the
\texttt{total\_amount\_due} accuracy ceiling on this corpus to
roughly 95\% with the current prompt.

\subsection{Construct validity}
\label{sec:eval-validity-construct}

Three documented cases of multiple-defensible-answer GT. On
\texttt{gb\_water\_affinitywater\_001} the GT \texttt{total\_amount\_due}
is \pounds46.00 (current charges); the bill also displays
\pounds235.14 (outstanding balance). On
\texttt{us\_gas\_peoplesgas\_001} the GT \texttt{account\_number}
is recorded with a particular whitespace pattern; the bill displays
the same digits with a different whitespace pattern. On
\texttt{us\_electricity\_publicservicecompany\_001} the GT
\texttt{provider\_name} is the legal entity ``Public Service Company
of Colorado''; the bill prominently displays ``Xcel Energy'', the
parent brand. The annotator picked one of two valid answers in
each case. Under the alternative defensible conventions the headline
numbers move by a few tenths of a percentage point and the modality
and provider rankings do not change.

\subsection{External validity and acknowledged limitations}
\label{sec:eval-validity-external}

Six bounds on external validity, each with its cause and the
specific finding it bounds.

\textbf{Corpus size.} $n=40$ active documents. Per-language and
per-utility-type slices are descriptive at this size; differences
of a few percentage points should not be over-read. The cause is
the manual collection constraint of
\S\ref{sec:eval-dataset-provenance}; expansion to $n\geq 200$
would lift the slice analyses to formally testable form.

\textbf{Language and script imbalance.} 20~English, 12~German,
7~Italian, 1~French; Latin script only. Non-Latin scripts (Arabic,
Cyrillic, CJK) are out of scope. The per-language results in
\S\ref{sec:eval-rq1-stratified} should not be extrapolated to
language families not exercised here.

\textbf{Single-annotator labelling.} Inter-annotator agreement
cannot be computed from a single labelling pass. The construct
threats above are the visible price; the multiple-defensible-answer
cases were caught by the author during error analysis but a second
pass would catch more.

\textbf{Single OCR engine.} The OCR-mediated pipeline uses
Tesseract only. Alternative OCR engines (Google Document AI,
PaddleOCR, AWS Textract) carry different error profiles, and the
``$\sim$50\% of OCR-mode failures are pure OCR'' finding is
therefore engine-specific. A drop-in replacement of the
\texttt{OCREngine} class (\S\ref{sec:design-pipeline}) is the
intended extension surface; running it remains future work.

\textbf{Two LLM providers, one model each.} Anthropic Sonnet 4.5
and OpenAI GPT-4o are the only two models tested. The
hallucination-vs-omission profile in \S\ref{sec:eval-rq3-orthogonal}
is therefore a two-point sample, not a population claim. Adding
Google Gemini, Mistral Large or open-weight models such as
Llama~3.2 Vision would broaden the provider axis and is the second
intended extension surface.

\textbf{Born-digital corpus only.} No scanned bills are included.
The OCR error ceiling identified in Lopresti's
analysis~\cite{lopresti2009ocrerrors} is therefore not exercised
at its hardest, and the OCR side of the modality comparison is
shown at its strongest. A follow-up pass with scanned reproductions
would isolate the scan-quality effect against the findings here.

% ------------------------------------------------------------
\section{Synthesis}
\label{sec:eval-synthesis}

Five findings carry forward to Chapter~\ref{chap:conclusion}.

\textbf{H1 is rejected on this corpus.} Vision-mediated extraction
did not match OCR-mediated extraction in aggregate, and the gap is
provider- and field-specific rather than uniform.

\textbf{H2 partially holds.} Vision is more expensive (consistent
with H2) and operationally simpler than OCR; vision is faster
end-to-end (against H2), and the latency reading is conditional on
single-process Tesseract.

\textbf{H3 is supported.} The within-provider modality gap differs
by provider, the providers occupy distinct points in
hallucination-versus-omission space, and the modality choice
interacts with provider choice rather than being orthogonal to it.

\textbf{Modality is a per-field decision.} Within-document modality
flipping (Figure~\ref{fig:eval-flips}) implies that no
single-modality system is uniformly best on this corpus, and the
empirical case for keeping both pipelines is the field-level
asymmetry, not aggregate accuracy.

\textbf{Half of OCR-mode failures are not Tesseract's fault.} The
four-bucket re-attribution shows that 45--49\% of OCR-mediated
failures are addressable by improving the LLM, the prompt or the
annotation rather than the preprocessing.

These five statements are the substrate Chapter~\ref{chap:conclusion}
turns into transferable claims and predictions for new domains.
.
%  Do NOT add \documentclass or \begin{document} here.
%  Follows Chapter 6 (Implementation); precedes Chapter 8
%  (Professional Issues / Conclusion).
%  Scope: tests every research question and pre-registered
%  hypothesis fixed in Chapter 4 against the run evidence written
%  by the system built in Chapter 6, including a comparison with
%  the alternative designs rejected in Chapter 5.
% ============================================================

\chapter{Evaluation}
\label{chap:evaluation}

This chapter asks whether the system specified in
Chapter~\ref{chap:spec} and built in Chapter~\ref{chap:impl} produces
the results predicted by the directional hypotheses fixed before any
run. Procedure and dataset are stated first
(\S\S\ref{sec:eval-procedure}--\ref{sec:eval-dataset}); each later
section reports against a target chosen before the data was seen.
Results are organised by research question
(\S\S\ref{sec:eval-headline}--\ref{sec:eval-rq3}), then by failure
mode and pipeline-design choice (\S\ref{sec:eval-failures}), then by
counterfactual against the alternatives discarded in
Chapter~\ref{chap:design} (\S\ref{sec:eval-alternatives}).
\S\ref{sec:eval-validity} states what these numbers cannot say;
\S\ref{sec:eval-synthesis} hands the transferable findings to
Chapter~\ref{chap:conclusion}.

% ------------------------------------------------------------
\section{Procedure and what counts as evidence}
\label{sec:eval-procedure}

Four runs were executed in April 2026, one per cell of the
$2\times 2$ factorial in Figure~\ref{fig:factorial}: Sonnet/OCR,
Sonnet/Vision, GPT-4o/OCR and GPT-4o/Vision. Each run processed all
forty documents in \texttt{data/dataset\_manifest.csv} with status
\texttt{active}, using the prompt template archived alongside the run
directory. Pipeline configuration, prompt and manifest snapshot were
written verbatim into each run directory at start time, so any later
run is a diff inspection of those three files away from the four
reported here.

The metric set is the one fixed in
Chapter~\ref{chap:spec}~\S\ref{sec:spec-evaluation}: field-level
accuracy with both-null cells excluded from the denominator,
document-level accuracy (a document counts as correct only if all
twelve fields are correct), and overall accuracy aggregated across
fields and documents. The four-class null taxonomy of
Figure~\ref{fig:null-taxonomy} is applied at field-comparison time
and persisted, so every later analysis below reads from the same
classification. Falsification follows the pre-registration: H1 is
rejected if vision is below OCR under the same provider, H2 is
rejected if any single condition is strictly cheaper, faster and
simpler than its peers, and H3 is rejected if the within-provider
modality gap is the same sign and magnitude across providers.

Three properties of the evidence are worth fixing now. Every run
completed end-to-end: 0/40 documents failed in any condition, so no
result below is restricted to a successful subset. No JSON-parse
failure occurred across the $40\times 12\times 4 = 1{,}920$ field
comparisons, so the parse-hardening of
Chapter~\ref{chap:design}~\S\ref{sec:design-execution} held and the
content findings are not contaminated by formatting failures. The
corpus is small ($n=40$) and errors cluster at the document level,
so per-language slices are reported as descriptive ranges rather than
as point estimates with confidence intervals. Significance testing is
not attempted, in line with the scope decision in
\S\ref{sec:spec-scope}.

% ------------------------------------------------------------
\section{Dataset and ground truth}
\label{sec:eval-dataset}

\subsection{Provenance}
\label{sec:eval-dataset-provenance}

Utility bills are personal-data instruments under GDPR
Article~4(1)~\cite{eu2016gdpr}. Random sampling from a billing
provider's archive, the textbook way to build a balanced
multilingual corpus, is therefore not feasible at
undergraduate-project scale. The corpus had to be assembled by hand.

Two streams contribute. Stream~A is bills released publicly by the
issuers themselves: redacted exemplars on customer-help pages,
sample bills on tariff-board portals, and reproductions on
billing-explainer guides. Stream~A accounts for the majority of the
corpus and covers the British, American, Australian, Canadian,
French, German and Italian subsets. Stream~B is bills contributed
by personal contacts under written consent for use in this study;
Stream~B sits inside the British and Italian subsets, fills two gaps
where public specimens were thin, and was used only after every
contributor signed a one-page consent form covering scope and
storage. Every document in either stream had personally identifiable
account, customer and address detail redacted before storage in the
repository, in line with the provenance controls
Chapter~\ref{chap:lsesp} sets out.

The corpus aims for within-budget diversity rather than statistical
representativeness, and the implications for external validity
appear in \S\ref{sec:eval-validity}. The provenance constraints are
not a defect of the dataset; they are the reason the dataset has
the shape it has.

\subsection{Composition}
\label{sec:eval-dataset-composition}

Forty documents are in scope. Twenty are English, twelve are German,
seven are Italian and one is French. Twenty-seven are electricity
bills, six are gas, and seven are water. Thirty-three unique issuers
appear; concentration is tracked in \S\ref{sec:eval-validity}
because four BC~Hydro bills, two Affinity~Water bills, two
Con~Edison bills, two Vattenfall Europe Sales bills, two Green
Planet Energy bills, two Lupatotina bills, and two PG\&E bills carry
greater weight on any per-issuer slice than a single document would.
Page-count distribution: eleven documents have at most two pages,
and twenty-nine documents (72.5\% of the corpus) include at least
one page beyond two. Figure~\ref{fig:eval-dataset-composition} shows
the three breakdowns on one row; the page-count panel returns in
\S\ref{sec:eval-failures-pagecap} when the two-page baseline is
defended empirically.

\begin{figure}[ht]
    \centering
    \includegraphics[width=\linewidth]{Chapters/figures/fig-7-7-dataset-composition.png}
    \caption[Dataset composition.]{The forty-document corpus by
    language, by utility type and by page count. The page-count panel
    marks the two-page pipeline cap; documents to the right of the
    dashed line have at least one page that neither pipeline
    processes. The cap is defended in
    \S\ref{sec:eval-failures-pagecap}.}
    \label{fig:eval-dataset-composition}
\end{figure}

\subsection{Annotation and the meaning of ``correct''}
\label{sec:eval-dataset-annotation}

A single annotator (the author) produced the per-document JSON files
in \texttt{data/ground\_truth/} using the format fixed in
\S\ref{sec:spec-schema}. Dates were recorded in ISO~8601, numerics
in locale-preserved form, and \texttt{null} was reserved for fields
genuinely absent from the bill rather than for fields that were
unreadable. The four-class taxonomy of
Figure~\ref{fig:null-taxonomy} therefore maps cleanly to predictions
without further annotation steps. Where a field admitted multiple
defensible answers the annotation records one of them, and the
choice is itself an annotation decision rather than a fact about the
bill; the specific cases are listed in \S\ref{sec:eval-validity} and
their effect on the headline numbers is bounded there.
Single-annotator coverage is acknowledged in
\S\ref{sec:eval-validity} as well, because inter-annotator agreement
cannot be computed from a single labelling pass.

% ------------------------------------------------------------
\section{Headline results}
\label{sec:eval-headline}

\begin{table}[ht]
    \centering
    \caption{Field-level overall accuracy and document-level accuracy
    across the four conditions. Both-null cells are excluded from the
    denominator (\S\ref{sec:spec-evaluation}). Source:
    \texttt{results/reports/four\_way\_20260429/comparison.json}.}
    \label{tab:eval-headline}
    \begin{tabular}{lcc}
        \hline
        Condition & Field-level overall & Document-level \\
        \hline
        Sonnet 4.5 / OCR-mediated     & $89.65\%$ & $45.0\%$ \\
        Sonnet 4.5 / Vision-mediated  & $87.55\%$ & $37.5\%$ \\
        GPT-4o / OCR-mediated         & $86.35\%$ & $30.0\%$ \\
        GPT-4o / Vision-mediated      & $85.33\%$ & $22.5\%$ \\
        \hline
    \end{tabular}
\end{table}

Table~\ref{tab:eval-headline} reports the four overall numbers and
the four document-level numbers. Figure~\ref{fig:eval-overall} is
the bar-chart view; Figure~\ref{fig:eval-heatmap} is the per-field
heatmap. At field level the four conditions cluster within a
4.32-percentage-point band; at document level the band widens to
22.5 points because a document is correct only when all twelve
fields are correct, and single-field misses across the schema
accumulate. The document-level signal is the more
production-relevant of the two
(Chapter~\ref{chap:spec}~\S\ref{sec:spec-evaluation}): a bill with
one wrong field cannot be ingested without manual review. None of
the four conditions clears 50\% document-level accuracy, which is
itself a finding about the gap between aggregate accuracy and
end-to-end usefulness on this document class at this prompt and
this schema.

\begin{figure}[ht]
    \centering
    \includegraphics[width=0.85\linewidth]{Chapters/figures/fig-7-1-overall-accuracy.png}
    \caption[Overall accuracy, four conditions.]{Field-level overall
    accuracy across the four conditions. Sonnet/OCR is the highest
    cell; GPT-4o/Vision is the lowest. The band is narrow at field
    level and wide at document level (Table~\ref{tab:eval-headline}).}
    \label{fig:eval-overall}
\end{figure}

\begin{figure}[ht]
    \centering
    \includegraphics[width=\linewidth]{Chapters/figures/fig-7-2-field-accuracy-heatmap.png}
    \caption[Per-field accuracy heatmap.]{Per-field accuracy across
    twelve fields and four conditions. \texttt{utility\_type},
    \texttt{currency} and \texttt{consumption\_unit} are uniformly
    high; \texttt{bill\_number}, \texttt{account\_number} and
    \texttt{provider\_name} are where the four conditions disagree
    the most. Source: \texttt{accuracy\_matrix} in
    \texttt{comparison.json}.}
    \label{fig:eval-heatmap}
\end{figure}

% ------------------------------------------------------------
\section{RQ1 --- does input modality matter?}
\label{sec:eval-rq1}

\subsection{Aggregate verdict on H1}
\label{sec:eval-rq1-aggregate}

H1 (vision $\geq$ OCR within provider) is rejected on this corpus.
For Sonnet, OCR-mediated extraction beats vision-mediated extraction
by $+2.10\,\mathrm{pp}$ at field level (89.65\% vs.\ 87.55\%); for
GPT-4o the gap is $+1.02\,\mathrm{pp}$ (86.35\% vs.\ 85.33\%). Both
gaps point the wrong way for H1. The strict reading is that H1 is
falsified for both providers. The cautious reading is that for
GPT-4o the gap is within the variation expected from $n=40$ and the
finding is then that the modality effect is much smaller for GPT-4o
than the literature
predicted~\cite{shen2026ocrornot}. The prior that
vision-mediated extraction would dominate OCR-mediated extraction
does not survive contact with this corpus.

\subsection{Where the gap actually lives}
\label{sec:eval-rq1-stratified}

The aggregate gap conceals slice-level differences large enough to
matter for system design. Figure~\ref{fig:eval-by-lang} stratifies
overall accuracy by language. The
``OCR\,$>$\,Vision'' finding is dragged by English: for Sonnet,
$\Delta_{\mathrm{en}} = +3.42\,\mathrm{pp}$; for GPT-4o,
$\Delta_{\mathrm{en}} = +2.15\,\mathrm{pp}$. On Italian
($n=7$) the GPT-4o ranking reverses to
$-2.72\,\mathrm{pp}$ (vision wins), and on French ($n=1$) Sonnet/Vision
is exactly correct on every field.
Figure~\ref{fig:eval-by-utility} is the same picture by utility
type: the smallest slice (water, $n=7$) shows the largest reversal
under Sonnet (vision narrowly behind on water but ahead on gas),
again from inside an aggregate that says vision is behind.

\begin{figure}[ht]
    \centering
    \includegraphics[width=0.95\linewidth]{Chapters/figures/fig-7-3-accuracy-by-language.png}
    \caption[Accuracy by language.]{Field-level accuracy by language
    for the four conditions. The aggregate ``OCR\,$>$\,Vision''
    finding is dragged by English; on Italian the GPT-4o ranking
    reverses, and on French ($n=1$) Sonnet/Vision is perfect.}
    \label{fig:eval-by-lang}
\end{figure}

\begin{figure}[ht]
    \centering
    \includegraphics[width=0.95\linewidth]{Chapters/figures/fig-7-4-accuracy-by-utility-type.png}
    \caption[Accuracy by utility type.]{Field-level accuracy by
    utility type. Water bills are the smallest slice ($n=7$) but
    show non-trivial modality reversals; the largest gas-side gap is
    GPT-4o, where OCR beats vision by $7.0\,\mathrm{pp}$.}
    \label{fig:eval-by-utility}
\end{figure}

\subsection{Field-level offsets that point in opposite directions}
\label{sec:eval-rq1-flips}

The aggregate gap is composed of per-field offsets that pull both
ways inside the same documents. Figure~\ref{fig:eval-flips} shows
the count of documents where one modality is correct on a given
field and the other is wrong, broken out per provider. For Sonnet,
\texttt{provider\_name} flips OCR-only-correct on six documents and
Vision-only-correct on two, a six-to-two asymmetry that contributes
the largest single chunk of the modality gap; on
\texttt{total\_amount\_due} the asymmetry inverts (Vision-only-correct
on four documents, OCR-only-correct on one). For GPT-4o,
\texttt{account\_number} flips Vision-only-correct on eleven
documents against OCR-only-correct on one, an asymmetry that pulls
GPT-4o/Vision \emph{ahead} on that field even though it is behind in
aggregate; \texttt{provider\_name} pulls in the opposite direction
(OCR-only-correct on nine documents against Vision-only-correct on
two).

\begin{figure}[ht]
    \centering
    \includegraphics[width=0.95\linewidth]{Chapters/figures/fig-7-11-modality-flip-counts.png}
    \caption[Modality-flip counts per field and provider.]{Per
    provider, per field, the count of documents where exactly one
    modality is correct. The two columns ``OCR-only correct'' and
    ``Vision-only correct'' sum to disagreement volume on that field
    for that provider, regardless of whether the modality wins or
    loses on that field overall.}
    \label{fig:eval-flips}
\end{figure}

The flip counts are the empirical case against any single-modality
system on this corpus. Twelve documents in the corpus carry at least
one such flip per provider, and on several documents the flips
oppose: a perfect modality router would beat the better of the two
single-modality systems by several percentage points
(\S\ref{sec:eval-alternatives}).

\subsection{Mechanisms behind the gap}
\label{sec:eval-rq1-mechanisms}

Four mechanisms account for most of the failures that drive the
aggregate gap. Each is identified by inspecting the wrong
predictions on representative documents in
\texttt{results/runs/}; named documents are given so the
attribution is auditable.

\textbf{(M1) Brand wordmark vs.\ legal name on
\texttt{provider\_name}.} Sonnet/Vision drops from 0.85 to 0.75 on
this field; GPT-4o drops from 0.825 to 0.65. Vision systematically
prefers the prominent brand wordmark on the bill header
(``ENGIE'', ``IREN'', ``Xcel Energy'') to the legal-entity surface
form recorded as ground truth. The OCR-mediated path sees the legal
suffix in the small print at the bottom of the bill; the
vision-mediated path anchors to whichever string is the largest
piece of type on the page. This is a header-salience effect, not a
legal-suffix-stripping effect: the ground truth records one of two
defensible surface forms and the modality systematically picks
different ones.

\textbf{(M2) Whitespace coercion on \texttt{account\_number}.} GPT-4o
in OCR mode scores 0.486 on this field while Sonnet in OCR mode
scores 0.750 and GPT-4o in vision mode scores 0.750. Inspection of
the GPT-4o/OCR misses shows GPT-4o stripping internal spaces from
account numbers it received with spaces in the OCR text, then
mismatching the spaced ground truth. Sonnet preserves the
whitespace; vision sees the layout directly and reproduces it. This
is a within-provider modality reversal: GPT-4o/Vision beats
GPT-4o/OCR on this field by 26.4~percentage points.

\textbf{(M3) Locale-aware separator handling on
\texttt{consumption\_amount}.} Both providers drop on this field
under vision: Sonnet from 0.923 to 0.846, GPT-4o from 0.923 to
0.821. The recurring failure pattern is European thousands
separators (``\texttt{16.250}'', ``\texttt{16\,250}'') being read by
the vision pipeline as a decimal (``\texttt{16.25}''). The
OCR-mediated path passes the literal string through the EU/US-aware
float parser of \S\ref{sec:design-normalisation}, which resolves
the locale at parse time using the per-document language flag.
Vision pipelines do their own number reading inside the model and
do not see that flag. The mechanism predicts that vision will
under-perform on European-locale numerics, which is the slice
pattern Figure~\ref{fig:eval-by-lang} shows.

\textbf{(M4) Field disambiguation on documents with multiple
candidates.} On bills that present several plausible totals
(``Current charges'', ``Outstanding balance'', ``Amount due''),
several plausible reference numbers (``Bill number'', ``Account
number'', ``Statement reference''), or several dates (``Bill date'',
``Issued'', ``Statement date''), both modalities sometimes pick the
wrong one. The clearest case is
\texttt{gb\_water\_affinitywater\_001}, where the ground-truth
\texttt{total\_amount\_due} is \pounds46.00 (current charges) but the
bill prominently displays \pounds235.14 (outstanding balance). The
disambiguation is field-construct ambiguity rather than a modality
problem: it appears in both modalities and is annotated as a
construct threat in \S\ref{sec:eval-validity}.

% ------------------------------------------------------------
\section{RQ2 --- the operational trade-off}
\label{sec:eval-rq2}

\subsection{Cost}
\label{sec:eval-rq2-cost}

H2 in its as-written form (vision more expensive) is supported.
Vision costs more per document for both providers: Sonnet/Vision
\$0.0153 against Sonnet/OCR \$0.0119 ($+28.8\%$), and GPT-4o/Vision
\$0.00938 against GPT-4o/OCR \$0.00828 ($+13.2\%$). Both pipelines
process the same first two pages of each bill (the design choice
defended in \S\ref{sec:eval-failures-pagecap}), so cost differences
trace directly to input-token volume: image tokens for the two
rasterised pages cost more than the text tokens for the
Tesseract-extracted text of the same two pages, on both providers'
2026 price lists. Figure~\ref{fig:eval-cost} is the per-document
view summed across the corpus.

\begin{figure}[ht]
    \centering
    \includegraphics[width=0.85\linewidth]{Chapters/figures/fig-7-5-cost-comparison.png}
    \caption[Per-document cost.]{Per-document estimated cost across
    the four conditions, summed over forty documents. Vision is
    more expensive than OCR for both providers; the gap is larger
    for Sonnet than for GPT-4o.}
    \label{fig:eval-cost}
\end{figure}

\subsection{Latency}
\label{sec:eval-rq2-latency}

H2 in its as-written form (vision slower) is rejected. Vision is
faster end-to-end for both providers. Mean wall-clock per document
is 9.65~seconds for GPT-4o/Vision against 14.44~seconds for
GPT-4o/OCR ($-33\%$), and 12.16~seconds for Sonnet/Vision against
16.35~seconds for Sonnet/OCR ($-26\%$). Figure~\ref{fig:eval-timing}
splits the time. Tesseract accounts for $\sim$8.7~seconds of every
OCR-mediated run; LLM-call latency rises in vision mode (Sonnet/Vision
8.6~seconds vs.\ Sonnet/OCR 3.7~seconds; GPT-4o/Vision 6.2~seconds
vs.\ GPT-4o/OCR 2.1~seconds) because the model now does the page
reading internally, but removing Tesseract is the larger effect on
total time on this hardware.

\begin{figure}[ht]
    \centering
    \includegraphics[width=0.85\linewidth]{Chapters/figures/fig-7-6-timing-breakdown.png}
    \caption[Timing breakdown.]{Mean per-document timings split into
    OCR and LLM stages. Tesseract dominates the OCR-mediated
    pipeline on this hardware; the LLM-call cost rises in vision
    mode but does not close the gap.}
    \label{fig:eval-timing}
\end{figure}

The latency ranking is hardware-conditioned. Tesseract was run
single-process on commodity CPU and is embarrassingly parallel, so
the OCR-side latency would shrink substantially on a multi-core
deployment. The latency ranking above is therefore a single-laptop
ranking, not a production-scale verdict; the production-scale
prediction is that the LLM-call dominates wall-clock under either
modality once OCR is parallelised, and the cross-modality latency
gap closes.

\subsection{Operational complexity}
\label{sec:eval-rq2-complexity}

Vision is simpler to deploy than OCR, and the ranking is
unambiguous. The OCR-mediated pipeline depends on Poppler (system
binary), Tesseract (system binary plus per-language data files for
\texttt{eng}, \texttt{deu}, \texttt{fra} and \texttt{ita}) and the
EU/US-aware float parser of
\S\ref{sec:design-normalisation}. The vision-mediated pipeline
depends on the provider SDK alone. Adding a fifth language to the
OCR-mediated pipeline requires a new Tesseract data file on every
host; adding a fifth language to the vision-mediated pipeline
requires nothing.

The H2 verdict is therefore mixed: vision is more expensive but
faster and operationally simpler. The trade-off direction H2
predicted is right; the per-axis directions split along provider
lines and along the cost-vs-latency axis in ways the as-written
hypothesis did not anticipate.

% ------------------------------------------------------------
\section{RQ3 --- does the choice of LLM matter?}
\label{sec:eval-rq3}

\subsection{Modality-by-provider interaction}
\label{sec:eval-rq3-interaction}

H3 (the modality gap differs by provider) is supported. The
within-provider modality gap is $+2.10\,\mathrm{pp}$ for Sonnet
and $+1.02\,\mathrm{pp}$ for GPT-4o, a gap-of-gaps of
$1.08\,\mathrm{pp}$. The interaction has a mechanistic reading:
Sonnet has more material to lose in vision mode because it is the
better OCR-mode reader on this corpus, and Sonnet's vision-mode
disadvantage concentrates on the M1 (header anchoring) and M3
(locale formatting) mechanisms above. GPT-4o is less prone to
header anchoring on the M1 cases inspected, partly because its
visible behaviour is to omit a field rather than substitute a
salient nearby string when the answer is uncertain
(\S\ref{sec:eval-rq3-orthogonal}).

\subsection{Provider effects orthogonal to modality}
\label{sec:eval-rq3-orthogonal}

Two provider effects appear in both modalities and are therefore
attributable to the model rather than to the input mode.

The first is currency contextual inference. The four BC~Hydro bills
have \texttt{currency} ground truth \texttt{CAD}, but the string
``CAD'' does not appear in the OCR text or on the page image: the
bills display the bare ``\$'' sign. Sonnet returns \texttt{CAD} on
all four under both OCR and vision; GPT-4o returns \texttt{USD} on
all four under both modalities. The signal Sonnet uses (``BC
Hydro'', ``Vancouver'', Canadian postal code) is available to both
providers; GPT-4o reads it and does not act on it. This single
mechanism accounts for the entire 0.90 vs.\ 1.00 currency gap in
Table~\ref{tab:eval-headline} between GPT-4o and Sonnet.

The second is the hallucination-versus-omission profile. Across
both modalities, Sonnet emits 20 hallucinations (8 OCR, 12 vision)
and 20 omissions (9 OCR, 11 vision); GPT-4o emits 5 hallucinations
(1 OCR, 4 vision) and 32 omissions (17 OCR, 15 vision). Sonnet
hallucinates roughly four times as often as GPT-4o; GPT-4o omits
where Sonnet guesses. Figure~\ref{fig:eval-pareto} plots the four
conditions on the two axes and shows that the providers occupy
distinct points, not different points on a single curve. Downstream
systems that punish wrong values should prefer GPT-4o; ones that
punish missing values should prefer Sonnet.

\begin{figure}[ht]
    \centering
    \includegraphics[width=0.78\linewidth]{Chapters/figures/fig-7-9-hallucination-omission-pareto.png}
    \caption[Hallucination-versus-omission Pareto.]{Per-condition
    counts of hallucinations and omissions across the corpus.
    Sonnet sits in the high-hallucination low-omission region under
    both modalities; GPT-4o sits in the low-hallucination
    higher-omission region. The two providers occupy distinct
    operating points, not different points on a single curve.}
    \label{fig:eval-pareto}
\end{figure}

\subsection{Cross-provider agreement and ensemble ceiling}
\label{sec:eval-rq3-agreement}

The two providers agree on the same answer for $\sim$80.4\% of OCR-mode
field comparisons and $\sim$78.1\% of vision-mode comparisons, where
``agree'' means both predictions normalise to the same value
regardless of correctness. Disagreements concentrate on the
mechanisms above: BC~Hydro \texttt{currency}, German credit-bill
totals, the Affinity Water \texttt{total\_amount\_due}
disambiguation, and \texttt{provider\_name} on bills with strong
brand wordmarks. A field-level vote ensemble of Sonnet and GPT-4o
in OCR mode would clear $\sim$91.4\% accuracy on this corpus
(against the better of the two cells alone, 89.65\% for Sonnet/OCR);
in vision mode the corresponding upper bound is $\sim$89.6\% against
87.55\% for Sonnet/Vision. The provider-level personality difference
H3 predicts is therefore exploitable, but the ceiling is not large
enough to outweigh the cost of running both providers per document.

% ------------------------------------------------------------
\section{Where failures actually come from}
\label{sec:eval-failures}

\subsection{The diagnosis heuristic and what it misses}
\label{sec:eval-failures-heuristic}

The diagnosis heuristic specified in
\S\ref{sec:design-evalframework-diagnosis} attributes a non-correct
OCR-mediated field to OCR if the ground-truth value is not a
case-insensitive substring of the stored OCR text, and to the LLM
otherwise. The heuristic is reproducible and inspectable, and it is
the right starting point. It also has known failure modes that
matter for the analysis below.

The first is contextual inference. On the BC~Hydro currency cases
above, the ground-truth string ``CAD'' is not in the OCR text, so
the strict heuristic attributes the GPT-4o failure to OCR. The
actual cause is that GPT-4o did not infer the currency from the
issuer name and address; Sonnet did infer it from exactly the same
OCR text. Attributing this to OCR is wrong in the sense that no OCR
change would have fixed it. The second is ground-truth phrasing: on
\texttt{us\_electricity\_publicservicecompany\_001} the GT
\texttt{provider\_name} is the legal name ``Public Service Company
of Colorado'', and the OCR text contains the brand wordmark ``Xcel
Energy'' rather than the legal name. The strict heuristic attributes
this to OCR; the actual cause is annotation choosing the legal name
where the brand wordmark would have been an equally defensible
answer. The third is locale-formatted numerics: on a bill that
displays ``\texttt{16.250}'' the GT records ``\texttt{16250}''
(thousands separator removed) and the heuristic again says OCR even
though the text was readable.

\subsection{Two-bucket vs.\ four-bucket re-attribution}
\label{sec:eval-failures-fourbucket}

Replacing two buckets with four addresses the failure modes above.
The four buckets are: \emph{Pure OCR} (Tesseract did not capture
the value), \emph{Pure LLM-extraction} (the LLM mis-read text that
was present), \emph{LLM-inference} (context for inference was
present and the model did not act on it; the BC~Hydro currency
cases sit here), and \emph{GT-ambiguity} (the annotation chose one
of multiple defensible answers; the PSCo and Affinity Water cases
sit here). The re-classification rule is: known multiple-answer
cases move to GT-ambiguity, currency and \texttt{utility\_type}
strict-OCR-failures move to LLM-inference, the rest stay where the
strict heuristic put them. Figure~\ref{fig:eval-attribution} shows
the result.

\begin{figure}[ht]
    \centering
    \includegraphics[width=0.95\linewidth]{Chapters/figures/fig-7-10-failure-attribution.png}
    \caption[Failure attribution: two-bucket vs.\ four-bucket.]{Failure
    attribution for the two OCR-mediated runs. Panel (a): the
    strict two-bucket heuristic of
    \S\ref{sec:design-evalframework-diagnosis}. Panel (b): four
    buckets after applying the re-classification rule of
    \S\ref{sec:eval-failures-fourbucket} to
    \texttt{derived/diagnosis\_audit.csv}.}
    \label{fig:eval-attribution}
\end{figure}

Of 47 Sonnet/OCR non-correct fields, 26 (55\%) are Pure OCR, 17
(36\%) are Pure LLM-extraction, 1 is LLM-inference and 3 are
GT-ambiguity. Of 61 GPT-4o/OCR non-correct fields, 31 (51\%) are
Pure OCR, 23 (38\%) are Pure LLM-extraction, 4 are LLM-inference
and 3 are GT-ambiguity. The headline finding survives in spirit: a
little over half of OCR-mode failures are addressable by improving
Tesseract; the remainder are addressable by improving the LLM, the
prompt or the annotation. The four-bucket reading sharpens that
finding by separating the cases where the OCR text was the problem
from the cases where the OCR text was fine and the model did not
act on what it had.

\subsection{Modality-specific patterns}
\label{sec:eval-failures-modality}

Two failure types appear under one input mode and not the other.
The OCR-mediated path hallucinates redacted spans into plausible
characters: on
\texttt{it\_electricity\_irenmercato\_001} the ground-truth
\texttt{bill\_number} is \texttt{502305934151}; the OCR text
contains the redacted form ``\texttt{50230XXXXXXX}'' and the LLM
returns a guessed completion in OCR mode while correctly returning
\texttt{null} (or the unredacted prefix) in vision mode where the
redaction box is visible as a layout element. The vision-mediated
path inverts the failure mode: on bills with prominent brand
wordmarks the model anchors to the wordmark on
\texttt{provider\_name} (M1), and on bills with European-locale
thousands separators the model collapses the separator (M3).

A third pattern affects every condition equally: sign handling on
credit and refund bills. On \texttt{de\_electricity\_naturstromhandel\_001}
the ground-truth \texttt{total\_amount\_due} is $-40.14$~EUR
because the bill is a refund (``Guthaben''). All four conditions
return $40.14$, dropping the sign. This is not a modality
strength on either side; it is a construct issue between the prompt
and the ground-truth convention, and it is recorded as a
limitation in \S\ref{sec:eval-validity}.

\subsection{The two-page baseline as a deliberate design choice}
\label{sec:eval-failures-pagecap}

Both pipelines process only the first two pages of each bill. The
cap is set explicitly in \texttt{src/pipeline.py} (lines 137 and
149: \texttt{ocr\_images = images[:2]} and
\texttt{vision\_images = images[:2]}). It is a deliberate baseline,
not a confound. On every bill in this corpus, every required field
that the bill carries appears within the first two pages: the
header (provider, utility type, bill number, dates), the totals
panel and the consumption box are universally on page one or two,
and later pages carry per-meter detail, terms and conditions,
multilingual translations or marketing inserts that the schema does
not target. The cap therefore costs almost nothing in accuracy
under this corpus while cutting input tokens by a factor of
$\sim$3.9 (the mean page count is 3.9 with the longest bill at
nine pages) and removing $\sim$5--7 seconds of Tesseract time on
the OCR side per long bill. Figure~\ref{fig:eval-page-cap}
visualises which documents are seen in full and which are
truncated.

\begin{figure}[ht]
    \centering
    \includegraphics[width=0.78\linewidth]{Chapters/figures/fig-7-8-page-cap-coverage.png}
    \caption[Two-page cap coverage.]{Per-document coverage under
    the two-page cap, applied to both pipelines. Dark segments are
    the pages either pipeline reads; light segments are the pages
    neither pipeline reads.}
    \label{fig:eval-page-cap}
\end{figure}

The cap is a property of the input distribution, not of the
algorithm. On a corpus that exercises later-page targets, changing
the slice from \texttt{[:2]} to a higher value, or to a function
that picks pages by a header-detection heuristic, would extend
coverage at proportional cost. The cost-vs-coverage curve for that
extension is the most direct empirical follow-up the corpus
suggests, and is left for Chapter~\ref{chap:conclusion}. The
relevant point here is that the cap is the reason cost and latency
are flat in document length on this corpus, and the reason the
modality-vs-modality comparison in \S\ref{sec:eval-rq1} is a
like-for-like comparison rather than a coverage-difference artefact.

% ------------------------------------------------------------
\section{Counterfactuals against the alternative designs}
\label{sec:eval-alternatives}

This section returns to the alternatives recorded in
Figure~\ref{fig:design-decisions} and tests them against the run
evidence rather than against intuition; the analysis closes the
A++ requirement that an evaluation chapter compare results expected
from theory for alternative designs with those obtained in
practice (Chapter~\ref{chap:spec}~\S\ref{sec:spec-evaluation}).

\subsection{Per-issuer templates}
\label{sec:eval-alternatives-templates}

The data ratifies the rejection of \S\ref{sec:design-decisions}.
Thirty-three unique issuers appear in the corpus and the field-level
error patterns are not positional: \texttt{provider\_name} fails on
brand-vs-legal selection (M1), \texttt{consumption\_amount} fails on
locale-aware separator handling (M3), and the
\texttt{total\_amount\_due} sign loss on credit bills is a semantic
question, not a coordinate-extraction question. None of those failure
modes is template-fixable; each one requires either an
issuer-specific rule or a flexible reading of the bill, and the
LLM-mediated approach gives the latter for free.

\subsection{Fine-tuned encoder--decoder (LayoutLMv3, Donut)}
\label{sec:eval-alternatives-finetune}

Forty documents is well below the floor at which Donut and
LayoutLMv3 are reported usable~\cite{kim2022donut,huang2022layoutlmv3};
provider concentration would inject training-distribution bias the
holdout cannot detect (BC~Hydro at four documents would be
disproportionately easy on any random split); and per-language
fine-tuning would be needed for the Italian and German subsets where
mechanisms M3 and the credit-bill sign issue occur. The LLM-mediated
approach pays a per-call cost in exchange for zero training data,
and at \$0.008--0.015 per document the per-call cost is below the
cost of one annotation hour at any reasonable rate.

\subsection{Provider ensembling}
\label{sec:eval-alternatives-ensemble}

The ensembling alternative is testable directly against the run
evidence (\S\ref{sec:eval-rq3-agreement}). A field-level vote of
Sonnet/OCR and GPT-4o/OCR clears $\sim$91.4\% on this corpus,
against 89.65\% and 86.35\% for the two cells alone. Most of the
gain is on the disagreement set (BC~Hydro \texttt{currency},
PSCo \texttt{provider\_name}, the German credit-bill totals where
Sonnet recovers the sign for one of the four cases). The ceiling
is real but moderate; pairing it with the cost ranking in
\S\ref{sec:eval-rq2-cost} (running two providers per document
roughly doubles the cost) puts the operating decision squarely in
the downstream-cost-tolerance domain.

\subsection{Single-modality system}
\label{sec:eval-alternatives-singlemod}

A vision-only or OCR-only system would have to cope with the
within-document modality flipping documented in
\S\ref{sec:eval-rq1-flips}. Twelve documents per provider have at
least one field where one modality is correct and the other wrong,
and several of those documents the directions oppose. A perfect
modality router would beat the better of the two single-modality
cells by roughly the same amount as a provider ensemble. The case
for keeping both pipelines is precisely this: not because either
wins outright, but because the conditions under which each is
preferable are field-, locale- and issuer-specific.

\subsection{Vision (or OCR) with all pages}
\label{sec:eval-alternatives-allpages}

The two-page cap is the design choice that makes both modalities
flat in document length on this corpus
(\S\ref{sec:eval-failures-pagecap}). Removing the cap on a corpus
where every required field already lives on the first two pages
would buy nothing in accuracy and cost roughly $3.9\times$ more in
input tokens (and proportionally more in OCR latency) on average.
Removing the cap on a different corpus where required fields appear
on later pages would buy accuracy at proportional cost, and the
cost-vs-coverage curve as a function of page count is the
empirical question Chapter~\ref{chap:conclusion} picks up. Either
way, the cap is the right choice for the corpus this study
evaluates.

% ------------------------------------------------------------
\section{What these numbers cannot say}
\label{sec:eval-validity}

\subsection{Internal validity}
\label{sec:eval-validity-internal}

Two internal threats matter. The diagnosis heuristic mis-attributes
the cases listed in \S\ref{sec:eval-failures-heuristic}; the
four-bucket re-classification of
\S\ref{sec:eval-failures-fourbucket} is the corrective. Sign
handling on credit-bill totals fails uniformly across all four
conditions because the prompt does not encode the
``\texttt{total\_amount\_due} is signed (negative on credit
bills)'' convention that the ground truth uses; this is a prompt
construct issue rather than a model failure, and it bounds the
\texttt{total\_amount\_due} accuracy ceiling on this corpus to
roughly 95\% with the current prompt.

\subsection{Construct validity}
\label{sec:eval-validity-construct}

Three documented cases of multiple-defensible-answer GT. On
\texttt{gb\_water\_affinitywater\_001} the GT \texttt{total\_amount\_due}
is \pounds46.00 (current charges); the bill also displays
\pounds235.14 (outstanding balance). On
\texttt{us\_gas\_peoplesgas\_001} the GT \texttt{account\_number}
is recorded with a particular whitespace pattern; the bill displays
the same digits with a different whitespace pattern. On
\texttt{us\_electricity\_publicservicecompany\_001} the GT
\texttt{provider\_name} is the legal entity ``Public Service Company
of Colorado''; the bill prominently displays ``Xcel Energy'', the
parent brand. The annotator picked one of two valid answers in
each case. Under the alternative defensible conventions the headline
numbers move by a few tenths of a percentage point and the modality
and provider rankings do not change.

\subsection{External validity and acknowledged limitations}
\label{sec:eval-validity-external}

Six bounds on external validity, each with its cause and the
specific finding it bounds.

\textbf{Corpus size.} $n=40$ active documents. Per-language and
per-utility-type slices are descriptive at this size; differences
of a few percentage points should not be over-read. The cause is
the manual collection constraint of
\S\ref{sec:eval-dataset-provenance}; expansion to $n\geq 200$
would lift the slice analyses to formally testable form.

\textbf{Language and script imbalance.} 20~English, 12~German,
7~Italian, 1~French; Latin script only. Non-Latin scripts (Arabic,
Cyrillic, CJK) are out of scope. The per-language results in
\S\ref{sec:eval-rq1-stratified} should not be extrapolated to
language families not exercised here.

\textbf{Single-annotator labelling.} Inter-annotator agreement
cannot be computed from a single labelling pass. The construct
threats above are the visible price; the multiple-defensible-answer
cases were caught by the author during error analysis but a second
pass would catch more.

\textbf{Single OCR engine.} The OCR-mediated pipeline uses
Tesseract only. Alternative OCR engines (Google Document AI,
PaddleOCR, AWS Textract) carry different error profiles, and the
``$\sim$50\% of OCR-mode failures are pure OCR'' finding is
therefore engine-specific. A drop-in replacement of the
\texttt{OCREngine} class (\S\ref{sec:design-pipeline}) is the
intended extension surface; running it remains future work.

\textbf{Two LLM providers, one model each.} Anthropic Sonnet 4.5
and OpenAI GPT-4o are the only two models tested. The
hallucination-vs-omission profile in \S\ref{sec:eval-rq3-orthogonal}
is therefore a two-point sample, not a population claim. Adding
Google Gemini, Mistral Large or open-weight models such as
Llama~3.2 Vision would broaden the provider axis and is the second
intended extension surface.

\textbf{Born-digital corpus only.} No scanned bills are included.
The OCR error ceiling identified in Lopresti's
analysis~\cite{lopresti2009ocrerrors} is therefore not exercised
at its hardest, and the OCR side of the modality comparison is
shown at its strongest. A follow-up pass with scanned reproductions
would isolate the scan-quality effect against the findings here.

% ------------------------------------------------------------
\section{Synthesis}
\label{sec:eval-synthesis}

Five findings carry forward to Chapter~\ref{chap:conclusion}.

\textbf{H1 is rejected on this corpus.} Vision-mediated extraction
did not match OCR-mediated extraction in aggregate, and the gap is
provider- and field-specific rather than uniform.

\textbf{H2 partially holds.} Vision is more expensive (consistent
with H2) and operationally simpler than OCR; vision is faster
end-to-end (against H2), and the latency reading is conditional on
single-process Tesseract.

\textbf{H3 is supported.} The within-provider modality gap differs
by provider, the providers occupy distinct points in
hallucination-versus-omission space, and the modality choice
interacts with provider choice rather than being orthogonal to it.

\textbf{Modality is a per-field decision.} Within-document modality
flipping (Figure~\ref{fig:eval-flips}) implies that no
single-modality system is uniformly best on this corpus, and the
empirical case for keeping both pipelines is the field-level
asymmetry, not aggregate accuracy.

\textbf{Half of OCR-mode failures are not Tesseract's fault.} The
four-bucket re-attribution shows that 45--49\% of OCR-mediated
failures are addressable by improving the LLM, the prompt or the
annotation rather than the preprocessing.

These five statements are the substrate Chapter~\ref{chap:conclusion}
turns into transferable claims and predictions for new domains.
.
%  Do NOT add \documentclass or \begin{document} here.
%  Follows Chapter 6 (Implementation); precedes Chapter 8
%  (Professional Issues / Conclusion).
%  Scope: tests every research question and pre-registered
%  hypothesis fixed in Chapter 4 against the run evidence written
%  by the system built in Chapter 6, including a comparison with
%  the alternative designs rejected in Chapter 5.
% ============================================================

\chapter{Evaluation}
\label{chap:evaluation}

This chapter asks whether the system specified in
Chapter~\ref{chap:spec} and built in Chapter~\ref{chap:impl} produces
the results predicted by the directional hypotheses fixed before any
run. Procedure and dataset are stated first
(\S\S\ref{sec:eval-procedure}--\ref{sec:eval-dataset}); each later
section reports against a target chosen before the data was seen.
Results are organised by research question
(\S\S\ref{sec:eval-headline}--\ref{sec:eval-rq3}), then by failure
mode and pipeline-design choice (\S\ref{sec:eval-failures}), then by
counterfactual against the alternatives discarded in
Chapter~\ref{chap:design} (\S\ref{sec:eval-alternatives}).
\S\ref{sec:eval-validity} states what these numbers cannot say;
\S\ref{sec:eval-synthesis} hands the transferable findings to
Chapter~\ref{chap:conclusion}.

% ------------------------------------------------------------
\section{Procedure and what counts as evidence}
\label{sec:eval-procedure}

Four runs were executed in April 2026, one per cell of the
$2\times 2$ factorial in Figure~\ref{fig:factorial}: Sonnet/OCR,
Sonnet/Vision, GPT-4o/OCR and GPT-4o/Vision. Each run processed all
forty documents in \texttt{data/dataset\_manifest.csv} with status
\texttt{active}, using the prompt template archived alongside the run
directory. Pipeline configuration, prompt and manifest snapshot were
written verbatim into each run directory at start time, so any later
run is a diff inspection of those three files away from the four
reported here.

The metric set is the one fixed in
Chapter~\ref{chap:spec}~\S\ref{sec:spec-evaluation}: field-level
accuracy with both-null cells excluded from the denominator,
document-level accuracy (a document counts as correct only if all
twelve fields are correct), and overall accuracy aggregated across
fields and documents. The four-class null taxonomy of
Figure~\ref{fig:null-taxonomy} is applied at field-comparison time
and persisted, so every later analysis below reads from the same
classification. Falsification follows the pre-registration: H1 is
rejected if vision is below OCR under the same provider, H2 is
rejected if any single condition is strictly cheaper, faster and
simpler than its peers, and H3 is rejected if the within-provider
modality gap is the same sign and magnitude across providers.

Three properties of the evidence are worth fixing now. Every run
completed end-to-end: 0/40 documents failed in any condition, so no
result below is restricted to a successful subset. No JSON-parse
failure occurred across the $40\times 12\times 4 = 1{,}920$ field
comparisons, so the parse-hardening of
Chapter~\ref{chap:design}~\S\ref{sec:design-execution} held and the
content findings are not contaminated by formatting failures. The
corpus is small ($n=40$) and errors cluster at the document level,
so per-language slices are reported as descriptive ranges rather than
as point estimates with confidence intervals. Significance testing is
not attempted, in line with the scope decision in
\S\ref{sec:spec-scope}.

% ------------------------------------------------------------
\section{Dataset and ground truth}
\label{sec:eval-dataset}

\subsection{Provenance}
\label{sec:eval-dataset-provenance}

Utility bills are personal-data instruments under GDPR
Article~4(1)~\cite{eu2016gdpr}. Random sampling from a billing
provider's archive, the textbook way to build a balanced
multilingual corpus, is therefore not feasible at
undergraduate-project scale. The corpus had to be assembled by hand.

Two streams contribute. Stream~A is bills released publicly by the
issuers themselves: redacted exemplars on customer-help pages,
sample bills on tariff-board portals, and reproductions on
billing-explainer guides. Stream~A accounts for the majority of the
corpus and covers the British, American, Australian, Canadian,
French, German and Italian subsets. Stream~B is bills contributed
by personal contacts under written consent for use in this study;
Stream~B sits inside the British and Italian subsets, fills two gaps
where public specimens were thin, and was used only after every
contributor signed a one-page consent form covering scope and
storage. Every document in either stream had personally identifiable
account, customer and address detail redacted before storage in the
repository, in line with the provenance controls
Chapter~\ref{chap:lsesp} sets out.

The corpus aims for within-budget diversity rather than statistical
representativeness, and the implications for external validity
appear in \S\ref{sec:eval-validity}. The provenance constraints are
not a defect of the dataset; they are the reason the dataset has
the shape it has.

\subsection{Composition}
\label{sec:eval-dataset-composition}

Forty documents are in scope. Twenty are English, twelve are German,
seven are Italian and one is French. Twenty-seven are electricity
bills, six are gas, and seven are water. Thirty-three unique issuers
appear; concentration is tracked in \S\ref{sec:eval-validity}
because four BC~Hydro bills, two Affinity~Water bills, two
Con~Edison bills, two Vattenfall Europe Sales bills, two Green
Planet Energy bills, two Lupatotina bills, and two PG\&E bills carry
greater weight on any per-issuer slice than a single document would.
Page-count distribution: eleven documents have at most two pages,
and twenty-nine documents (72.5\% of the corpus) include at least
one page beyond two. Figure~\ref{fig:eval-dataset-composition} shows
the three breakdowns on one row; the page-count panel returns in
\S\ref{sec:eval-failures-pagecap} when the two-page baseline is
defended empirically.

\begin{figure}[ht]
    \centering
    \includegraphics[width=\linewidth]{Chapters/figures/fig-7-7-dataset-composition.png}
    \caption[Dataset composition.]{The forty-document corpus by
    language, by utility type and by page count. The page-count panel
    marks the two-page pipeline cap; documents to the right of the
    dashed line have at least one page that neither pipeline
    processes. The cap is defended in
    \S\ref{sec:eval-failures-pagecap}.}
    \label{fig:eval-dataset-composition}
\end{figure}

\subsection{Annotation and the meaning of ``correct''}
\label{sec:eval-dataset-annotation}

A single annotator (the author) produced the per-document JSON files
in \texttt{data/ground\_truth/} using the format fixed in
\S\ref{sec:spec-schema}. Dates were recorded in ISO~8601, numerics
in locale-preserved form, and \texttt{null} was reserved for fields
genuinely absent from the bill rather than for fields that were
unreadable. The four-class taxonomy of
Figure~\ref{fig:null-taxonomy} therefore maps cleanly to predictions
without further annotation steps. Where a field admitted multiple
defensible answers the annotation records one of them, and the
choice is itself an annotation decision rather than a fact about the
bill; the specific cases are listed in \S\ref{sec:eval-validity} and
their effect on the headline numbers is bounded there.
Single-annotator coverage is acknowledged in
\S\ref{sec:eval-validity} as well, because inter-annotator agreement
cannot be computed from a single labelling pass.

% ------------------------------------------------------------
\section{Headline results}
\label{sec:eval-headline}

\begin{table}[ht]
    \centering
    \caption{Field-level overall accuracy and document-level accuracy
    across the four conditions. Both-null cells are excluded from the
    denominator (\S\ref{sec:spec-evaluation}). Source:
    \texttt{results/reports/four\_way\_20260429/comparison.json}.}
    \label{tab:eval-headline}
    \begin{tabular}{lcc}
        \hline
        Condition & Field-level overall & Document-level \\
        \hline
        Sonnet 4.5 / OCR-mediated     & $89.65\%$ & $45.0\%$ \\
        Sonnet 4.5 / Vision-mediated  & $87.55\%$ & $37.5\%$ \\
        GPT-4o / OCR-mediated         & $86.35\%$ & $30.0\%$ \\
        GPT-4o / Vision-mediated      & $85.33\%$ & $22.5\%$ \\
        \hline
    \end{tabular}
\end{table}

Table~\ref{tab:eval-headline} reports the four overall numbers and
the four document-level numbers. Figure~\ref{fig:eval-overall} is
the bar-chart view; Figure~\ref{fig:eval-heatmap} is the per-field
heatmap. At field level the four conditions cluster within a
4.32-percentage-point band; at document level the band widens to
22.5 points because a document is correct only when all twelve
fields are correct, and single-field misses across the schema
accumulate. The document-level signal is the more
production-relevant of the two
(Chapter~\ref{chap:spec}~\S\ref{sec:spec-evaluation}): a bill with
one wrong field cannot be ingested without manual review. None of
the four conditions clears 50\% document-level accuracy, which is
itself a finding about the gap between aggregate accuracy and
end-to-end usefulness on this document class at this prompt and
this schema.

\begin{figure}[ht]
    \centering
    \includegraphics[width=0.85\linewidth]{Chapters/figures/fig-7-1-overall-accuracy.png}
    \caption[Overall accuracy, four conditions.]{Field-level overall
    accuracy across the four conditions. Sonnet/OCR is the highest
    cell; GPT-4o/Vision is the lowest. The band is narrow at field
    level and wide at document level (Table~\ref{tab:eval-headline}).}
    \label{fig:eval-overall}
\end{figure}

\begin{figure}[ht]
    \centering
    \includegraphics[width=\linewidth]{Chapters/figures/fig-7-2-field-accuracy-heatmap.png}
    \caption[Per-field accuracy heatmap.]{Per-field accuracy across
    twelve fields and four conditions. \texttt{utility\_type},
    \texttt{currency} and \texttt{consumption\_unit} are uniformly
    high; \texttt{bill\_number}, \texttt{account\_number} and
    \texttt{provider\_name} are where the four conditions disagree
    the most. Source: \texttt{accuracy\_matrix} in
    \texttt{comparison.json}.}
    \label{fig:eval-heatmap}
\end{figure}

% ------------------------------------------------------------
\section{RQ1 --- does input modality matter?}
\label{sec:eval-rq1}

\subsection{Aggregate verdict on H1}
\label{sec:eval-rq1-aggregate}

H1 (vision $\geq$ OCR within provider) is rejected on this corpus.
For Sonnet, OCR-mediated extraction beats vision-mediated extraction
by $+2.10\,\mathrm{pp}$ at field level (89.65\% vs.\ 87.55\%); for
GPT-4o the gap is $+1.02\,\mathrm{pp}$ (86.35\% vs.\ 85.33\%). Both
gaps point the wrong way for H1. The strict reading is that H1 is
falsified for both providers. The cautious reading is that for
GPT-4o the gap is within the variation expected from $n=40$ and the
finding is then that the modality effect is much smaller for GPT-4o
than the literature
predicted~\cite{shen2026ocrornot}. The prior that
vision-mediated extraction would dominate OCR-mediated extraction
does not survive contact with this corpus.

\subsection{Where the gap actually lives}
\label{sec:eval-rq1-stratified}

The aggregate gap conceals slice-level differences large enough to
matter for system design. Figure~\ref{fig:eval-by-lang} stratifies
overall accuracy by language. The
``OCR\,$>$\,Vision'' finding is dragged by English: for Sonnet,
$\Delta_{\mathrm{en}} = +3.42\,\mathrm{pp}$; for GPT-4o,
$\Delta_{\mathrm{en}} = +2.15\,\mathrm{pp}$. On Italian
($n=7$) the GPT-4o ranking reverses to
$-2.72\,\mathrm{pp}$ (vision wins), and on French ($n=1$) Sonnet/Vision
is exactly correct on every field.
Figure~\ref{fig:eval-by-utility} is the same picture by utility
type: the smallest slice (water, $n=7$) shows the largest reversal
under Sonnet (vision narrowly behind on water but ahead on gas),
again from inside an aggregate that says vision is behind.

\begin{figure}[ht]
    \centering
    \includegraphics[width=0.95\linewidth]{Chapters/figures/fig-7-3-accuracy-by-language.png}
    \caption[Accuracy by language.]{Field-level accuracy by language
    for the four conditions. The aggregate ``OCR\,$>$\,Vision''
    finding is dragged by English; on Italian the GPT-4o ranking
    reverses, and on French ($n=1$) Sonnet/Vision is perfect.}
    \label{fig:eval-by-lang}
\end{figure}

\begin{figure}[ht]
    \centering
    \includegraphics[width=0.95\linewidth]{Chapters/figures/fig-7-4-accuracy-by-utility-type.png}
    \caption[Accuracy by utility type.]{Field-level accuracy by
    utility type. Water bills are the smallest slice ($n=7$) but
    show non-trivial modality reversals; the largest gas-side gap is
    GPT-4o, where OCR beats vision by $7.0\,\mathrm{pp}$.}
    \label{fig:eval-by-utility}
\end{figure}

\subsection{Field-level offsets that point in opposite directions}
\label{sec:eval-rq1-flips}

The aggregate gap is composed of per-field offsets that pull both
ways inside the same documents. Figure~\ref{fig:eval-flips} shows
the count of documents where one modality is correct on a given
field and the other is wrong, broken out per provider. For Sonnet,
\texttt{provider\_name} flips OCR-only-correct on six documents and
Vision-only-correct on two, a six-to-two asymmetry that contributes
the largest single chunk of the modality gap; on
\texttt{total\_amount\_due} the asymmetry inverts (Vision-only-correct
on four documents, OCR-only-correct on one). For GPT-4o,
\texttt{account\_number} flips Vision-only-correct on eleven
documents against OCR-only-correct on one, an asymmetry that pulls
GPT-4o/Vision \emph{ahead} on that field even though it is behind in
aggregate; \texttt{provider\_name} pulls in the opposite direction
(OCR-only-correct on nine documents against Vision-only-correct on
two).

\begin{figure}[ht]
    \centering
    \includegraphics[width=0.95\linewidth]{Chapters/figures/fig-7-11-modality-flip-counts.png}
    \caption[Modality-flip counts per field and provider.]{Per
    provider, per field, the count of documents where exactly one
    modality is correct. The two columns ``OCR-only correct'' and
    ``Vision-only correct'' sum to disagreement volume on that field
    for that provider, regardless of whether the modality wins or
    loses on that field overall.}
    \label{fig:eval-flips}
\end{figure}

The flip counts are the empirical case against any single-modality
system on this corpus. Twelve documents in the corpus carry at least
one such flip per provider, and on several documents the flips
oppose: a perfect modality router would beat the better of the two
single-modality systems by several percentage points
(\S\ref{sec:eval-alternatives}).

\subsection{Mechanisms behind the gap}
\label{sec:eval-rq1-mechanisms}

Four mechanisms account for most of the failures that drive the
aggregate gap. Each is identified by inspecting the wrong
predictions on representative documents in
\texttt{results/runs/}; named documents are given so the
attribution is auditable.

\textbf{(M1) Brand wordmark vs.\ legal name on
\texttt{provider\_name}.} Sonnet/Vision drops from 0.85 to 0.75 on
this field; GPT-4o drops from 0.825 to 0.65. Vision systematically
prefers the prominent brand wordmark on the bill header
(``ENGIE'', ``IREN'', ``Xcel Energy'') to the legal-entity surface
form recorded as ground truth. The OCR-mediated path sees the legal
suffix in the small print at the bottom of the bill; the
vision-mediated path anchors to whichever string is the largest
piece of type on the page. This is a header-salience effect, not a
legal-suffix-stripping effect: the ground truth records one of two
defensible surface forms and the modality systematically picks
different ones.

\textbf{(M2) Whitespace coercion on \texttt{account\_number}.} GPT-4o
in OCR mode scores 0.486 on this field while Sonnet in OCR mode
scores 0.750 and GPT-4o in vision mode scores 0.750. Inspection of
the GPT-4o/OCR misses shows GPT-4o stripping internal spaces from
account numbers it received with spaces in the OCR text, then
mismatching the spaced ground truth. Sonnet preserves the
whitespace; vision sees the layout directly and reproduces it. This
is a within-provider modality reversal: GPT-4o/Vision beats
GPT-4o/OCR on this field by 26.4~percentage points.

\textbf{(M3) Locale-aware separator handling on
\texttt{consumption\_amount}.} Both providers drop on this field
under vision: Sonnet from 0.923 to 0.846, GPT-4o from 0.923 to
0.821. The recurring failure pattern is European thousands
separators (``\texttt{16.250}'', ``\texttt{16\,250}'') being read by
the vision pipeline as a decimal (``\texttt{16.25}''). The
OCR-mediated path passes the literal string through the EU/US-aware
float parser of \S\ref{sec:design-normalisation}, which resolves
the locale at parse time using the per-document language flag.
Vision pipelines do their own number reading inside the model and
do not see that flag. The mechanism predicts that vision will
under-perform on European-locale numerics, which is the slice
pattern Figure~\ref{fig:eval-by-lang} shows.

\textbf{(M4) Field disambiguation on documents with multiple
candidates.} On bills that present several plausible totals
(``Current charges'', ``Outstanding balance'', ``Amount due''),
several plausible reference numbers (``Bill number'', ``Account
number'', ``Statement reference''), or several dates (``Bill date'',
``Issued'', ``Statement date''), both modalities sometimes pick the
wrong one. The clearest case is
\texttt{gb\_water\_affinitywater\_001}, where the ground-truth
\texttt{total\_amount\_due} is \pounds46.00 (current charges) but the
bill prominently displays \pounds235.14 (outstanding balance). The
disambiguation is field-construct ambiguity rather than a modality
problem: it appears in both modalities and is annotated as a
construct threat in \S\ref{sec:eval-validity}.

% ------------------------------------------------------------
\section{RQ2 --- the operational trade-off}
\label{sec:eval-rq2}

\subsection{Cost}
\label{sec:eval-rq2-cost}

H2 in its as-written form (vision more expensive) is supported.
Vision costs more per document for both providers: Sonnet/Vision
\$0.0153 against Sonnet/OCR \$0.0119 ($+28.8\%$), and GPT-4o/Vision
\$0.00938 against GPT-4o/OCR \$0.00828 ($+13.2\%$). Both pipelines
process the same first two pages of each bill (the design choice
defended in \S\ref{sec:eval-failures-pagecap}), so cost differences
trace directly to input-token volume: image tokens for the two
rasterised pages cost more than the text tokens for the
Tesseract-extracted text of the same two pages, on both providers'
2026 price lists. Figure~\ref{fig:eval-cost} is the per-document
view summed across the corpus.

\begin{figure}[ht]
    \centering
    \includegraphics[width=0.85\linewidth]{Chapters/figures/fig-7-5-cost-comparison.png}
    \caption[Per-document cost.]{Per-document estimated cost across
    the four conditions, summed over forty documents. Vision is
    more expensive than OCR for both providers; the gap is larger
    for Sonnet than for GPT-4o.}
    \label{fig:eval-cost}
\end{figure}

\subsection{Latency}
\label{sec:eval-rq2-latency}

H2 in its as-written form (vision slower) is rejected. Vision is
faster end-to-end for both providers. Mean wall-clock per document
is 9.65~seconds for GPT-4o/Vision against 14.44~seconds for
GPT-4o/OCR ($-33\%$), and 12.16~seconds for Sonnet/Vision against
16.35~seconds for Sonnet/OCR ($-26\%$). Figure~\ref{fig:eval-timing}
splits the time. Tesseract accounts for $\sim$8.7~seconds of every
OCR-mediated run; LLM-call latency rises in vision mode (Sonnet/Vision
8.6~seconds vs.\ Sonnet/OCR 3.7~seconds; GPT-4o/Vision 6.2~seconds
vs.\ GPT-4o/OCR 2.1~seconds) because the model now does the page
reading internally, but removing Tesseract is the larger effect on
total time on this hardware.

\begin{figure}[ht]
    \centering
    \includegraphics[width=0.85\linewidth]{Chapters/figures/fig-7-6-timing-breakdown.png}
    \caption[Timing breakdown.]{Mean per-document timings split into
    OCR and LLM stages. Tesseract dominates the OCR-mediated
    pipeline on this hardware; the LLM-call cost rises in vision
    mode but does not close the gap.}
    \label{fig:eval-timing}
\end{figure}

The latency ranking is hardware-conditioned. Tesseract was run
single-process on commodity CPU and is embarrassingly parallel, so
the OCR-side latency would shrink substantially on a multi-core
deployment. The latency ranking above is therefore a single-laptop
ranking, not a production-scale verdict; the production-scale
prediction is that the LLM-call dominates wall-clock under either
modality once OCR is parallelised, and the cross-modality latency
gap closes.

\subsection{Operational complexity}
\label{sec:eval-rq2-complexity}

Vision is simpler to deploy than OCR, and the ranking is
unambiguous. The OCR-mediated pipeline depends on Poppler (system
binary), Tesseract (system binary plus per-language data files for
\texttt{eng}, \texttt{deu}, \texttt{fra} and \texttt{ita}) and the
EU/US-aware float parser of
\S\ref{sec:design-normalisation}. The vision-mediated pipeline
depends on the provider SDK alone. Adding a fifth language to the
OCR-mediated pipeline requires a new Tesseract data file on every
host; adding a fifth language to the vision-mediated pipeline
requires nothing.

The H2 verdict is therefore mixed: vision is more expensive but
faster and operationally simpler. The trade-off direction H2
predicted is right; the per-axis directions split along provider
lines and along the cost-vs-latency axis in ways the as-written
hypothesis did not anticipate.

% ------------------------------------------------------------
\section{RQ3 --- does the choice of LLM matter?}
\label{sec:eval-rq3}

\subsection{Modality-by-provider interaction}
\label{sec:eval-rq3-interaction}

H3 (the modality gap differs by provider) is supported. The
within-provider modality gap is $+2.10\,\mathrm{pp}$ for Sonnet
and $+1.02\,\mathrm{pp}$ for GPT-4o, a gap-of-gaps of
$1.08\,\mathrm{pp}$. The interaction has a mechanistic reading:
Sonnet has more material to lose in vision mode because it is the
better OCR-mode reader on this corpus, and Sonnet's vision-mode
disadvantage concentrates on the M1 (header anchoring) and M3
(locale formatting) mechanisms above. GPT-4o is less prone to
header anchoring on the M1 cases inspected, partly because its
visible behaviour is to omit a field rather than substitute a
salient nearby string when the answer is uncertain
(\S\ref{sec:eval-rq3-orthogonal}).

\subsection{Provider effects orthogonal to modality}
\label{sec:eval-rq3-orthogonal}

Two provider effects appear in both modalities and are therefore
attributable to the model rather than to the input mode.

The first is currency contextual inference. The four BC~Hydro bills
have \texttt{currency} ground truth \texttt{CAD}, but the string
``CAD'' does not appear in the OCR text or on the page image: the
bills display the bare ``\$'' sign. Sonnet returns \texttt{CAD} on
all four under both OCR and vision; GPT-4o returns \texttt{USD} on
all four under both modalities. The signal Sonnet uses (``BC
Hydro'', ``Vancouver'', Canadian postal code) is available to both
providers; GPT-4o reads it and does not act on it. This single
mechanism accounts for the entire 0.90 vs.\ 1.00 currency gap in
Table~\ref{tab:eval-headline} between GPT-4o and Sonnet.

The second is the hallucination-versus-omission profile. Across
both modalities, Sonnet emits 20 hallucinations (8 OCR, 12 vision)
and 20 omissions (9 OCR, 11 vision); GPT-4o emits 5 hallucinations
(1 OCR, 4 vision) and 32 omissions (17 OCR, 15 vision). Sonnet
hallucinates roughly four times as often as GPT-4o; GPT-4o omits
where Sonnet guesses. Figure~\ref{fig:eval-pareto} plots the four
conditions on the two axes and shows that the providers occupy
distinct points, not different points on a single curve. Downstream
systems that punish wrong values should prefer GPT-4o; ones that
punish missing values should prefer Sonnet.

\begin{figure}[ht]
    \centering
    \includegraphics[width=0.78\linewidth]{Chapters/figures/fig-7-9-hallucination-omission-pareto.png}
    \caption[Hallucination-versus-omission Pareto.]{Per-condition
    counts of hallucinations and omissions across the corpus.
    Sonnet sits in the high-hallucination low-omission region under
    both modalities; GPT-4o sits in the low-hallucination
    higher-omission region. The two providers occupy distinct
    operating points, not different points on a single curve.}
    \label{fig:eval-pareto}
\end{figure}

\subsection{Cross-provider agreement and ensemble ceiling}
\label{sec:eval-rq3-agreement}

The two providers agree on the same answer for $\sim$80.4\% of OCR-mode
field comparisons and $\sim$78.1\% of vision-mode comparisons, where
``agree'' means both predictions normalise to the same value
regardless of correctness. Disagreements concentrate on the
mechanisms above: BC~Hydro \texttt{currency}, German credit-bill
totals, the Affinity Water \texttt{total\_amount\_due}
disambiguation, and \texttt{provider\_name} on bills with strong
brand wordmarks. A field-level vote ensemble of Sonnet and GPT-4o
in OCR mode would clear $\sim$91.4\% accuracy on this corpus
(against the better of the two cells alone, 89.65\% for Sonnet/OCR);
in vision mode the corresponding upper bound is $\sim$89.6\% against
87.55\% for Sonnet/Vision. The provider-level personality difference
H3 predicts is therefore exploitable, but the ceiling is not large
enough to outweigh the cost of running both providers per document.

% ------------------------------------------------------------
\section{Where failures actually come from}
\label{sec:eval-failures}

\subsection{The diagnosis heuristic and what it misses}
\label{sec:eval-failures-heuristic}

The diagnosis heuristic specified in
\S\ref{sec:design-evalframework-diagnosis} attributes a non-correct
OCR-mediated field to OCR if the ground-truth value is not a
case-insensitive substring of the stored OCR text, and to the LLM
otherwise. The heuristic is reproducible and inspectable, and it is
the right starting point. It also has known failure modes that
matter for the analysis below.

The first is contextual inference. On the BC~Hydro currency cases
above, the ground-truth string ``CAD'' is not in the OCR text, so
the strict heuristic attributes the GPT-4o failure to OCR. The
actual cause is that GPT-4o did not infer the currency from the
issuer name and address; Sonnet did infer it from exactly the same
OCR text. Attributing this to OCR is wrong in the sense that no OCR
change would have fixed it. The second is ground-truth phrasing: on
\texttt{us\_electricity\_publicservicecompany\_001} the GT
\texttt{provider\_name} is the legal name ``Public Service Company
of Colorado'', and the OCR text contains the brand wordmark ``Xcel
Energy'' rather than the legal name. The strict heuristic attributes
this to OCR; the actual cause is annotation choosing the legal name
where the brand wordmark would have been an equally defensible
answer. The third is locale-formatted numerics: on a bill that
displays ``\texttt{16.250}'' the GT records ``\texttt{16250}''
(thousands separator removed) and the heuristic again says OCR even
though the text was readable.

\subsection{Two-bucket vs.\ four-bucket re-attribution}
\label{sec:eval-failures-fourbucket}

Replacing two buckets with four addresses the failure modes above.
The four buckets are: \emph{Pure OCR} (Tesseract did not capture
the value), \emph{Pure LLM-extraction} (the LLM mis-read text that
was present), \emph{LLM-inference} (context for inference was
present and the model did not act on it; the BC~Hydro currency
cases sit here), and \emph{GT-ambiguity} (the annotation chose one
of multiple defensible answers; the PSCo and Affinity Water cases
sit here). The re-classification rule is: known multiple-answer
cases move to GT-ambiguity, currency and \texttt{utility\_type}
strict-OCR-failures move to LLM-inference, the rest stay where the
strict heuristic put them. Figure~\ref{fig:eval-attribution} shows
the result.

\begin{figure}[ht]
    \centering
    \includegraphics[width=0.95\linewidth]{Chapters/figures/fig-7-10-failure-attribution.png}
    \caption[Failure attribution: two-bucket vs.\ four-bucket.]{Failure
    attribution for the two OCR-mediated runs. Panel (a): the
    strict two-bucket heuristic of
    \S\ref{sec:design-evalframework-diagnosis}. Panel (b): four
    buckets after applying the re-classification rule of
    \S\ref{sec:eval-failures-fourbucket} to
    \texttt{derived/diagnosis\_audit.csv}.}
    \label{fig:eval-attribution}
\end{figure}

Of 47 Sonnet/OCR non-correct fields, 26 (55\%) are Pure OCR, 17
(36\%) are Pure LLM-extraction, 1 is LLM-inference and 3 are
GT-ambiguity. Of 61 GPT-4o/OCR non-correct fields, 31 (51\%) are
Pure OCR, 23 (38\%) are Pure LLM-extraction, 4 are LLM-inference
and 3 are GT-ambiguity. The headline finding survives in spirit: a
little over half of OCR-mode failures are addressable by improving
Tesseract; the remainder are addressable by improving the LLM, the
prompt or the annotation. The four-bucket reading sharpens that
finding by separating the cases where the OCR text was the problem
from the cases where the OCR text was fine and the model did not
act on what it had.

\subsection{Modality-specific patterns}
\label{sec:eval-failures-modality}

Two failure types appear under one input mode and not the other.
The OCR-mediated path hallucinates redacted spans into plausible
characters: on
\texttt{it\_electricity\_irenmercato\_001} the ground-truth
\texttt{bill\_number} is \texttt{502305934151}; the OCR text
contains the redacted form ``\texttt{50230XXXXXXX}'' and the LLM
returns a guessed completion in OCR mode while correctly returning
\texttt{null} (or the unredacted prefix) in vision mode where the
redaction box is visible as a layout element. The vision-mediated
path inverts the failure mode: on bills with prominent brand
wordmarks the model anchors to the wordmark on
\texttt{provider\_name} (M1), and on bills with European-locale
thousands separators the model collapses the separator (M3).

A third pattern affects every condition equally: sign handling on
credit and refund bills. On \texttt{de\_electricity\_naturstromhandel\_001}
the ground-truth \texttt{total\_amount\_due} is $-40.14$~EUR
because the bill is a refund (``Guthaben''). All four conditions
return $40.14$, dropping the sign. This is not a modality
strength on either side; it is a construct issue between the prompt
and the ground-truth convention, and it is recorded as a
limitation in \S\ref{sec:eval-validity}.

\subsection{The two-page baseline as a deliberate design choice}
\label{sec:eval-failures-pagecap}

Both pipelines process only the first two pages of each bill. The
cap is set explicitly in \texttt{src/pipeline.py} (lines 137 and
149: \texttt{ocr\_images = images[:2]} and
\texttt{vision\_images = images[:2]}). It is a deliberate baseline,
not a confound. On every bill in this corpus, every required field
that the bill carries appears within the first two pages: the
header (provider, utility type, bill number, dates), the totals
panel and the consumption box are universally on page one or two,
and later pages carry per-meter detail, terms and conditions,
multilingual translations or marketing inserts that the schema does
not target. The cap therefore costs almost nothing in accuracy
under this corpus while cutting input tokens by a factor of
$\sim$3.9 (the mean page count is 3.9 with the longest bill at
nine pages) and removing $\sim$5--7 seconds of Tesseract time on
the OCR side per long bill. Figure~\ref{fig:eval-page-cap}
visualises which documents are seen in full and which are
truncated.

\begin{figure}[ht]
    \centering
    \includegraphics[width=0.78\linewidth]{Chapters/figures/fig-7-8-page-cap-coverage.png}
    \caption[Two-page cap coverage.]{Per-document coverage under
    the two-page cap, applied to both pipelines. Dark segments are
    the pages either pipeline reads; light segments are the pages
    neither pipeline reads.}
    \label{fig:eval-page-cap}
\end{figure}

The cap is a property of the input distribution, not of the
algorithm. On a corpus that exercises later-page targets, changing
the slice from \texttt{[:2]} to a higher value, or to a function
that picks pages by a header-detection heuristic, would extend
coverage at proportional cost. The cost-vs-coverage curve for that
extension is the most direct empirical follow-up the corpus
suggests, and is left for Chapter~\ref{chap:conclusion}. The
relevant point here is that the cap is the reason cost and latency
are flat in document length on this corpus, and the reason the
modality-vs-modality comparison in \S\ref{sec:eval-rq1} is a
like-for-like comparison rather than a coverage-difference artefact.

% ------------------------------------------------------------
\section{Counterfactuals against the alternative designs}
\label{sec:eval-alternatives}

This section returns to the alternatives recorded in
Figure~\ref{fig:design-decisions} and tests them against the run
evidence rather than against intuition; the analysis closes the
A++ requirement that an evaluation chapter compare results expected
from theory for alternative designs with those obtained in
practice (Chapter~\ref{chap:spec}~\S\ref{sec:spec-evaluation}).

\subsection{Per-issuer templates}
\label{sec:eval-alternatives-templates}

The data ratifies the rejection of \S\ref{sec:design-decisions}.
Thirty-three unique issuers appear in the corpus and the field-level
error patterns are not positional: \texttt{provider\_name} fails on
brand-vs-legal selection (M1), \texttt{consumption\_amount} fails on
locale-aware separator handling (M3), and the
\texttt{total\_amount\_due} sign loss on credit bills is a semantic
question, not a coordinate-extraction question. None of those failure
modes is template-fixable; each one requires either an
issuer-specific rule or a flexible reading of the bill, and the
LLM-mediated approach gives the latter for free.

\subsection{Fine-tuned encoder--decoder (LayoutLMv3, Donut)}
\label{sec:eval-alternatives-finetune}

Forty documents is well below the floor at which Donut and
LayoutLMv3 are reported usable~\cite{kim2022donut,huang2022layoutlmv3};
provider concentration would inject training-distribution bias the
holdout cannot detect (BC~Hydro at four documents would be
disproportionately easy on any random split); and per-language
fine-tuning would be needed for the Italian and German subsets where
mechanisms M3 and the credit-bill sign issue occur. The LLM-mediated
approach pays a per-call cost in exchange for zero training data,
and at \$0.008--0.015 per document the per-call cost is below the
cost of one annotation hour at any reasonable rate.

\subsection{Provider ensembling}
\label{sec:eval-alternatives-ensemble}

The ensembling alternative is testable directly against the run
evidence (\S\ref{sec:eval-rq3-agreement}). A field-level vote of
Sonnet/OCR and GPT-4o/OCR clears $\sim$91.4\% on this corpus,
against 89.65\% and 86.35\% for the two cells alone. Most of the
gain is on the disagreement set (BC~Hydro \texttt{currency},
PSCo \texttt{provider\_name}, the German credit-bill totals where
Sonnet recovers the sign for one of the four cases). The ceiling
is real but moderate; pairing it with the cost ranking in
\S\ref{sec:eval-rq2-cost} (running two providers per document
roughly doubles the cost) puts the operating decision squarely in
the downstream-cost-tolerance domain.

\subsection{Single-modality system}
\label{sec:eval-alternatives-singlemod}

A vision-only or OCR-only system would have to cope with the
within-document modality flipping documented in
\S\ref{sec:eval-rq1-flips}. Twelve documents per provider have at
least one field where one modality is correct and the other wrong,
and several of those documents the directions oppose. A perfect
modality router would beat the better of the two single-modality
cells by roughly the same amount as a provider ensemble. The case
for keeping both pipelines is precisely this: not because either
wins outright, but because the conditions under which each is
preferable are field-, locale- and issuer-specific.

\subsection{Vision (or OCR) with all pages}
\label{sec:eval-alternatives-allpages}

The two-page cap is the design choice that makes both modalities
flat in document length on this corpus
(\S\ref{sec:eval-failures-pagecap}). Removing the cap on a corpus
where every required field already lives on the first two pages
would buy nothing in accuracy and cost roughly $3.9\times$ more in
input tokens (and proportionally more in OCR latency) on average.
Removing the cap on a different corpus where required fields appear
on later pages would buy accuracy at proportional cost, and the
cost-vs-coverage curve as a function of page count is the
empirical question Chapter~\ref{chap:conclusion} picks up. Either
way, the cap is the right choice for the corpus this study
evaluates.

% ------------------------------------------------------------
\section{What these numbers cannot say}
\label{sec:eval-validity}

\subsection{Internal validity}
\label{sec:eval-validity-internal}

Two internal threats matter. The diagnosis heuristic mis-attributes
the cases listed in \S\ref{sec:eval-failures-heuristic}; the
four-bucket re-classification of
\S\ref{sec:eval-failures-fourbucket} is the corrective. Sign
handling on credit-bill totals fails uniformly across all four
conditions because the prompt does not encode the
``\texttt{total\_amount\_due} is signed (negative on credit
bills)'' convention that the ground truth uses; this is a prompt
construct issue rather than a model failure, and it bounds the
\texttt{total\_amount\_due} accuracy ceiling on this corpus to
roughly 95\% with the current prompt.

\subsection{Construct validity}
\label{sec:eval-validity-construct}

Three documented cases of multiple-defensible-answer GT. On
\texttt{gb\_water\_affinitywater\_001} the GT \texttt{total\_amount\_due}
is \pounds46.00 (current charges); the bill also displays
\pounds235.14 (outstanding balance). On
\texttt{us\_gas\_peoplesgas\_001} the GT \texttt{account\_number}
is recorded with a particular whitespace pattern; the bill displays
the same digits with a different whitespace pattern. On
\texttt{us\_electricity\_publicservicecompany\_001} the GT
\texttt{provider\_name} is the legal entity ``Public Service Company
of Colorado''; the bill prominently displays ``Xcel Energy'', the
parent brand. The annotator picked one of two valid answers in
each case. Under the alternative defensible conventions the headline
numbers move by a few tenths of a percentage point and the modality
and provider rankings do not change.

\subsection{External validity and acknowledged limitations}
\label{sec:eval-validity-external}

Six bounds on external validity, each with its cause and the
specific finding it bounds.

\textbf{Corpus size.} $n=40$ active documents. Per-language and
per-utility-type slices are descriptive at this size; differences
of a few percentage points should not be over-read. The cause is
the manual collection constraint of
\S\ref{sec:eval-dataset-provenance}; expansion to $n\geq 200$
would lift the slice analyses to formally testable form.

\textbf{Language and script imbalance.} 20~English, 12~German,
7~Italian, 1~French; Latin script only. Non-Latin scripts (Arabic,
Cyrillic, CJK) are out of scope. The per-language results in
\S\ref{sec:eval-rq1-stratified} should not be extrapolated to
language families not exercised here.

\textbf{Single-annotator labelling.} Inter-annotator agreement
cannot be computed from a single labelling pass. The construct
threats above are the visible price; the multiple-defensible-answer
cases were caught by the author during error analysis but a second
pass would catch more.

\textbf{Single OCR engine.} The OCR-mediated pipeline uses
Tesseract only. Alternative OCR engines (Google Document AI,
PaddleOCR, AWS Textract) carry different error profiles, and the
``$\sim$50\% of OCR-mode failures are pure OCR'' finding is
therefore engine-specific. A drop-in replacement of the
\texttt{OCREngine} class (\S\ref{sec:design-pipeline}) is the
intended extension surface; running it remains future work.

\textbf{Two LLM providers, one model each.} Anthropic Sonnet 4.5
and OpenAI GPT-4o are the only two models tested. The
hallucination-vs-omission profile in \S\ref{sec:eval-rq3-orthogonal}
is therefore a two-point sample, not a population claim. Adding
Google Gemini, Mistral Large or open-weight models such as
Llama~3.2 Vision would broaden the provider axis and is the second
intended extension surface.

\textbf{Born-digital corpus only.} No scanned bills are included.
The OCR error ceiling identified in Lopresti's
analysis~\cite{lopresti2009ocrerrors} is therefore not exercised
at its hardest, and the OCR side of the modality comparison is
shown at its strongest. A follow-up pass with scanned reproductions
would isolate the scan-quality effect against the findings here.

% ------------------------------------------------------------
\section{Synthesis}
\label{sec:eval-synthesis}

Five findings carry forward to Chapter~\ref{chap:conclusion}.

\textbf{H1 is rejected on this corpus.} Vision-mediated extraction
did not match OCR-mediated extraction in aggregate, and the gap is
provider- and field-specific rather than uniform.

\textbf{H2 partially holds.} Vision is more expensive (consistent
with H2) and operationally simpler than OCR; vision is faster
end-to-end (against H2), and the latency reading is conditional on
single-process Tesseract.

\textbf{H3 is supported.} The within-provider modality gap differs
by provider, the providers occupy distinct points in
hallucination-versus-omission space, and the modality choice
interacts with provider choice rather than being orthogonal to it.

\textbf{Modality is a per-field decision.} Within-document modality
flipping (Figure~\ref{fig:eval-flips}) implies that no
single-modality system is uniformly best on this corpus, and the
empirical case for keeping both pipelines is the field-level
asymmetry, not aggregate accuracy.

\textbf{Half of OCR-mode failures are not Tesseract's fault.} The
four-bucket re-attribution shows that 45--49\% of OCR-mediated
failures are addressable by improving the LLM, the prompt or the
annotation rather than the preprocessing.

These five statements are the substrate Chapter~\ref{chap:conclusion}
turns into transferable claims and predictions for new domains.
.
%  Do NOT add \documentclass or \begin{document} here.
%  Follows Chapter 6 (Implementation); precedes Chapter 8
%  (Professional Issues / Conclusion).
%  Scope: tests every research question and pre-registered
%  hypothesis fixed in Chapter 4 against the run evidence written
%  by the system built in Chapter 6, including a comparison with
%  the alternative designs rejected in Chapter 5.
% ============================================================

\chapter{Evaluation}
\label{chap:evaluation}

This chapter asks whether the system specified in
Chapter~\ref{chap:spec} and built in Chapter~\ref{chap:impl} produces
the results predicted by the directional hypotheses fixed before any
run. Procedure and dataset are stated first
(\S\S\ref{sec:eval-procedure}--\ref{sec:eval-dataset}); each later
section reports against a target chosen before the data was seen.
Results are organised by research question
(\S\S\ref{sec:eval-headline}--\ref{sec:eval-rq3}), then by failure
mode and pipeline-design choice (\S\ref{sec:eval-failures}), then by
counterfactual against the alternatives discarded in
Chapter~\ref{chap:design} (\S\ref{sec:eval-alternatives}).
\S\ref{sec:eval-validity} states what these numbers cannot say;
\S\ref{sec:eval-synthesis} hands the transferable findings to
Chapter~\ref{chap:conclusion}.

% ------------------------------------------------------------
\section{Procedure and what counts as evidence}
\label{sec:eval-procedure}

Four runs were executed in April 2026, one per cell of the
$2\times 2$ factorial in Figure~\ref{fig:factorial}: Sonnet/OCR,
Sonnet/Vision, GPT-4o/OCR and GPT-4o/Vision. Each run processed all
forty documents in \texttt{data/dataset\_manifest.csv} with status
\texttt{active}, using the prompt template archived alongside the run
directory. Pipeline configuration, prompt and manifest snapshot were
written verbatim into each run directory at start time, so any later
run is a diff inspection of those three files away from the four
reported here.

The metric set is the one fixed in
Chapter~\ref{chap:spec}~\S\ref{sec:spec-evaluation}: field-level
accuracy with both-null cells excluded from the denominator,
document-level accuracy (a document counts as correct only if all
twelve fields are correct), and overall accuracy aggregated across
fields and documents. The four-class null taxonomy of
Figure~\ref{fig:null-taxonomy} is applied at field-comparison time
and persisted, so every later analysis below reads from the same
classification. Falsification follows the pre-registration: H1 is
rejected if vision is below OCR under the same provider, H2 is
rejected if any single condition is strictly cheaper, faster and
simpler than its peers, and H3 is rejected if the within-provider
modality gap is the same sign and magnitude across providers.

Three properties of the evidence are worth fixing now. Every run
completed end-to-end: 0/40 documents failed in any condition, so no
result below is restricted to a successful subset. No JSON-parse
failure occurred across the $40\times 12\times 4 = 1{,}920$ field
comparisons, so the parse-hardening of
Chapter~\ref{chap:design}~\S\ref{sec:design-execution} held and the
content findings are not contaminated by formatting failures. The
corpus is small ($n=40$) and errors cluster at the document level,
so per-language slices are reported as descriptive ranges rather than
as point estimates with confidence intervals. Significance testing is
not attempted, in line with the scope decision in
\S\ref{sec:spec-scope}.

% ------------------------------------------------------------
\section{Dataset and ground truth}
\label{sec:eval-dataset}

\subsection{Provenance}
\label{sec:eval-dataset-provenance}

Utility bills are personal-data instruments under GDPR
Article~4(1)~\cite{eu2016gdpr}. Random sampling from a billing
provider's archive, the textbook way to build a balanced
multilingual corpus, is therefore not feasible at
undergraduate-project scale. The corpus had to be assembled by hand.

Two streams contribute. Stream~A is bills released publicly by the
issuers themselves: redacted exemplars on customer-help pages,
sample bills on tariff-board portals, and reproductions on
billing-explainer guides. Stream~A accounts for the majority of the
corpus and covers the British, American, Australian, Canadian,
French, German and Italian subsets. Stream~B is bills contributed
by personal contacts under written consent for use in this study;
Stream~B sits inside the British and Italian subsets, fills two gaps
where public specimens were thin, and was used only after every
contributor signed a one-page consent form covering scope and
storage. Every document in either stream had personally identifiable
account, customer and address detail redacted before storage in the
repository, in line with the provenance controls
Chapter~\ref{chap:lsesp} sets out.

The corpus aims for within-budget diversity rather than statistical
representativeness, and the implications for external validity
appear in \S\ref{sec:eval-validity}. The provenance constraints are
not a defect of the dataset; they are the reason the dataset has
the shape it has.

\subsection{Composition}
\label{sec:eval-dataset-composition}

Forty documents are in scope. Twenty are English, twelve are German,
seven are Italian and one is French. Twenty-seven are electricity
bills, six are gas, and seven are water. Thirty-three unique issuers
appear; concentration is tracked in \S\ref{sec:eval-validity}
because four BC~Hydro bills, two Affinity~Water bills, two
Con~Edison bills, two Vattenfall Europe Sales bills, two Green
Planet Energy bills, two Lupatotina bills, and two PG\&E bills carry
greater weight on any per-issuer slice than a single document would.
Page-count distribution: eleven documents have at most two pages,
and twenty-nine documents (72.5\% of the corpus) include at least
one page beyond two. Figure~\ref{fig:eval-dataset-composition} shows
the three breakdowns on one row; the page-count panel returns in
\S\ref{sec:eval-failures-pagecap} when the two-page baseline is
defended empirically.

\begin{figure}[ht]
    \centering
    \includegraphics[width=\linewidth]{Chapters/figures/fig-7-7-dataset-composition.png}
    \caption[Dataset composition.]{The forty-document corpus by
    language, by utility type and by page count. The page-count panel
    marks the two-page pipeline cap; documents to the right of the
    dashed line have at least one page that neither pipeline
    processes. The cap is defended in
    \S\ref{sec:eval-failures-pagecap}.}
    \label{fig:eval-dataset-composition}
\end{figure}

\subsection{Annotation and the meaning of ``correct''}
\label{sec:eval-dataset-annotation}

A single annotator (the author) produced the per-document JSON files
in \texttt{data/ground\_truth/} using the format fixed in
\S\ref{sec:spec-schema}. Dates were recorded in ISO~8601, numerics
in locale-preserved form, and \texttt{null} was reserved for fields
genuinely absent from the bill rather than for fields that were
unreadable. The four-class taxonomy of
Figure~\ref{fig:null-taxonomy} therefore maps cleanly to predictions
without further annotation steps. Where a field admitted multiple
defensible answers the annotation records one of them, and the
choice is itself an annotation decision rather than a fact about the
bill; the specific cases are listed in \S\ref{sec:eval-validity} and
their effect on the headline numbers is bounded there.
Single-annotator coverage is acknowledged in
\S\ref{sec:eval-validity} as well, because inter-annotator agreement
cannot be computed from a single labelling pass.

% ------------------------------------------------------------
\section{Headline results}
\label{sec:eval-headline}

\begin{table}[ht]
    \centering
    \caption{Field-level overall accuracy and document-level accuracy
    across the four conditions. Both-null cells are excluded from the
    denominator (\S\ref{sec:spec-evaluation}). Source:
    \texttt{results/reports/four\_way\_20260429/comparison.json}.}
    \label{tab:eval-headline}
    \begin{tabular}{lcc}
        \hline
        Condition & Field-level overall & Document-level \\
        \hline
        Sonnet 4.5 / OCR-mediated     & $89.65\%$ & $45.0\%$ \\
        Sonnet 4.5 / Vision-mediated  & $87.55\%$ & $37.5\%$ \\
        GPT-4o / OCR-mediated         & $86.35\%$ & $30.0\%$ \\
        GPT-4o / Vision-mediated      & $85.33\%$ & $22.5\%$ \\
        \hline
    \end{tabular}
\end{table}

Table~\ref{tab:eval-headline} reports the four overall numbers and
the four document-level numbers. Figure~\ref{fig:eval-overall} is
the bar-chart view; Figure~\ref{fig:eval-heatmap} is the per-field
heatmap. At field level the four conditions cluster within a
4.32-percentage-point band; at document level the band widens to
22.5 points because a document is correct only when all twelve
fields are correct, and single-field misses across the schema
accumulate. The document-level signal is the more
production-relevant of the two
(Chapter~\ref{chap:spec}~\S\ref{sec:spec-evaluation}): a bill with
one wrong field cannot be ingested without manual review. None of
the four conditions clears 50\% document-level accuracy, which is
itself a finding about the gap between aggregate accuracy and
end-to-end usefulness on this document class at this prompt and
this schema.

\begin{figure}[ht]
    \centering
    \includegraphics[width=0.85\linewidth]{Chapters/figures/fig-7-1-overall-accuracy.png}
    \caption[Overall accuracy, four conditions.]{Field-level overall
    accuracy across the four conditions. Sonnet/OCR is the highest
    cell; GPT-4o/Vision is the lowest. The band is narrow at field
    level and wide at document level (Table~\ref{tab:eval-headline}).}
    \label{fig:eval-overall}
\end{figure}

\begin{figure}[ht]
    \centering
    \includegraphics[width=\linewidth]{Chapters/figures/fig-7-2-field-accuracy-heatmap.png}
    \caption[Per-field accuracy heatmap.]{Per-field accuracy across
    twelve fields and four conditions. \texttt{utility\_type},
    \texttt{currency} and \texttt{consumption\_unit} are uniformly
    high; \texttt{bill\_number}, \texttt{account\_number} and
    \texttt{provider\_name} are where the four conditions disagree
    the most. Source: \texttt{accuracy\_matrix} in
    \texttt{comparison.json}.}
    \label{fig:eval-heatmap}
\end{figure}

% ------------------------------------------------------------
\section{RQ1 --- does input modality matter?}
\label{sec:eval-rq1}

\subsection{Aggregate verdict on H1}
\label{sec:eval-rq1-aggregate}

H1 (vision $\geq$ OCR within provider) is rejected on this corpus.
For Sonnet, OCR-mediated extraction beats vision-mediated extraction
by $+2.10\,\mathrm{pp}$ at field level (89.65\% vs.\ 87.55\%); for
GPT-4o the gap is $+1.02\,\mathrm{pp}$ (86.35\% vs.\ 85.33\%). Both
gaps point the wrong way for H1. The strict reading is that H1 is
falsified for both providers. The cautious reading is that for
GPT-4o the gap is within the variation expected from $n=40$ and the
finding is then that the modality effect is much smaller for GPT-4o
than the literature
predicted~\cite{shen2026ocrornot}. The prior that
vision-mediated extraction would dominate OCR-mediated extraction
does not survive contact with this corpus.

\subsection{Where the gap actually lives}
\label{sec:eval-rq1-stratified}

The aggregate gap conceals slice-level differences large enough to
matter for system design. Figure~\ref{fig:eval-by-lang} stratifies
overall accuracy by language. The
``OCR\,$>$\,Vision'' finding is dragged by English: for Sonnet,
$\Delta_{\mathrm{en}} = +3.42\,\mathrm{pp}$; for GPT-4o,
$\Delta_{\mathrm{en}} = +2.15\,\mathrm{pp}$. On Italian
($n=7$) the GPT-4o ranking reverses to
$-2.72\,\mathrm{pp}$ (vision wins), and on French ($n=1$) Sonnet/Vision
is exactly correct on every field.
Figure~\ref{fig:eval-by-utility} is the same picture by utility
type: the smallest slice (water, $n=7$) shows the largest reversal
under Sonnet (vision narrowly behind on water but ahead on gas),
again from inside an aggregate that says vision is behind.

\begin{figure}[ht]
    \centering
    \includegraphics[width=0.95\linewidth]{Chapters/figures/fig-7-3-accuracy-by-language.png}
    \caption[Accuracy by language.]{Field-level accuracy by language
    for the four conditions. The aggregate ``OCR\,$>$\,Vision''
    finding is dragged by English; on Italian the GPT-4o ranking
    reverses, and on French ($n=1$) Sonnet/Vision is perfect.}
    \label{fig:eval-by-lang}
\end{figure}

\begin{figure}[ht]
    \centering
    \includegraphics[width=0.95\linewidth]{Chapters/figures/fig-7-4-accuracy-by-utility-type.png}
    \caption[Accuracy by utility type.]{Field-level accuracy by
    utility type. Water bills are the smallest slice ($n=7$) but
    show non-trivial modality reversals; the largest gas-side gap is
    GPT-4o, where OCR beats vision by $7.0\,\mathrm{pp}$.}
    \label{fig:eval-by-utility}
\end{figure}

\subsection{Field-level offsets that point in opposite directions}
\label{sec:eval-rq1-flips}

The aggregate gap is composed of per-field offsets that pull both
ways inside the same documents. Figure~\ref{fig:eval-flips} shows
the count of documents where one modality is correct on a given
field and the other is wrong, broken out per provider. For Sonnet,
\texttt{provider\_name} flips OCR-only-correct on six documents and
Vision-only-correct on two, a six-to-two asymmetry that contributes
the largest single chunk of the modality gap; on
\texttt{total\_amount\_due} the asymmetry inverts (Vision-only-correct
on four documents, OCR-only-correct on one). For GPT-4o,
\texttt{account\_number} flips Vision-only-correct on eleven
documents against OCR-only-correct on one, an asymmetry that pulls
GPT-4o/Vision \emph{ahead} on that field even though it is behind in
aggregate; \texttt{provider\_name} pulls in the opposite direction
(OCR-only-correct on nine documents against Vision-only-correct on
two).

\begin{figure}[ht]
    \centering
    \includegraphics[width=0.95\linewidth]{Chapters/figures/fig-7-11-modality-flip-counts.png}
    \caption[Modality-flip counts per field and provider.]{Per
    provider, per field, the count of documents where exactly one
    modality is correct. The two columns ``OCR-only correct'' and
    ``Vision-only correct'' sum to disagreement volume on that field
    for that provider, regardless of whether the modality wins or
    loses on that field overall.}
    \label{fig:eval-flips}
\end{figure}

The flip counts are the empirical case against any single-modality
system on this corpus. Twelve documents in the corpus carry at least
one such flip per provider, and on several documents the flips
oppose: a perfect modality router would beat the better of the two
single-modality systems by several percentage points
(\S\ref{sec:eval-alternatives}).

\subsection{Mechanisms behind the gap}
\label{sec:eval-rq1-mechanisms}

Four mechanisms account for most of the failures that drive the
aggregate gap. Each is identified by inspecting the wrong
predictions on representative documents in
\texttt{results/runs/}; named documents are given so the
attribution is auditable.

\textbf{(M1) Brand wordmark vs.\ legal name on
\texttt{provider\_name}.} Sonnet/Vision drops from 0.85 to 0.75 on
this field; GPT-4o drops from 0.825 to 0.65. Vision systematically
prefers the prominent brand wordmark on the bill header
(``ENGIE'', ``IREN'', ``Xcel Energy'') to the legal-entity surface
form recorded as ground truth. The OCR-mediated path sees the legal
suffix in the small print at the bottom of the bill; the
vision-mediated path anchors to whichever string is the largest
piece of type on the page. This is a header-salience effect, not a
legal-suffix-stripping effect: the ground truth records one of two
defensible surface forms and the modality systematically picks
different ones.

\textbf{(M2) Whitespace coercion on \texttt{account\_number}.} GPT-4o
in OCR mode scores 0.486 on this field while Sonnet in OCR mode
scores 0.750 and GPT-4o in vision mode scores 0.750. Inspection of
the GPT-4o/OCR misses shows GPT-4o stripping internal spaces from
account numbers it received with spaces in the OCR text, then
mismatching the spaced ground truth. Sonnet preserves the
whitespace; vision sees the layout directly and reproduces it. This
is a within-provider modality reversal: GPT-4o/Vision beats
GPT-4o/OCR on this field by 26.4~percentage points.

\textbf{(M3) Locale-aware separator handling on
\texttt{consumption\_amount}.} Both providers drop on this field
under vision: Sonnet from 0.923 to 0.846, GPT-4o from 0.923 to
0.821. The recurring failure pattern is European thousands
separators (``\texttt{16.250}'', ``\texttt{16\,250}'') being read by
the vision pipeline as a decimal (``\texttt{16.25}''). The
OCR-mediated path passes the literal string through the EU/US-aware
float parser of \S\ref{sec:design-normalisation}, which resolves
the locale at parse time using the per-document language flag.
Vision pipelines do their own number reading inside the model and
do not see that flag. The mechanism predicts that vision will
under-perform on European-locale numerics, which is the slice
pattern Figure~\ref{fig:eval-by-lang} shows.

\textbf{(M4) Field disambiguation on documents with multiple
candidates.} On bills that present several plausible totals
(``Current charges'', ``Outstanding balance'', ``Amount due''),
several plausible reference numbers (``Bill number'', ``Account
number'', ``Statement reference''), or several dates (``Bill date'',
``Issued'', ``Statement date''), both modalities sometimes pick the
wrong one. The clearest case is
\texttt{gb\_water\_affinitywater\_001}, where the ground-truth
\texttt{total\_amount\_due} is \pounds46.00 (current charges) but the
bill prominently displays \pounds235.14 (outstanding balance). The
disambiguation is field-construct ambiguity rather than a modality
problem: it appears in both modalities and is annotated as a
construct threat in \S\ref{sec:eval-validity}.

% ------------------------------------------------------------
\section{RQ2 --- the operational trade-off}
\label{sec:eval-rq2}

\subsection{Cost}
\label{sec:eval-rq2-cost}

H2 in its as-written form (vision more expensive) is supported.
Vision costs more per document for both providers: Sonnet/Vision
\$0.0153 against Sonnet/OCR \$0.0119 ($+28.8\%$), and GPT-4o/Vision
\$0.00938 against GPT-4o/OCR \$0.00828 ($+13.2\%$). Both pipelines
process the same first two pages of each bill (the design choice
defended in \S\ref{sec:eval-failures-pagecap}), so cost differences
trace directly to input-token volume: image tokens for the two
rasterised pages cost more than the text tokens for the
Tesseract-extracted text of the same two pages, on both providers'
2026 price lists. Figure~\ref{fig:eval-cost} is the per-document
view summed across the corpus.

\begin{figure}[ht]
    \centering
    \includegraphics[width=0.85\linewidth]{Chapters/figures/fig-7-5-cost-comparison.png}
    \caption[Per-document cost.]{Per-document estimated cost across
    the four conditions, summed over forty documents. Vision is
    more expensive than OCR for both providers; the gap is larger
    for Sonnet than for GPT-4o.}
    \label{fig:eval-cost}
\end{figure}

\subsection{Latency}
\label{sec:eval-rq2-latency}

H2 in its as-written form (vision slower) is rejected. Vision is
faster end-to-end for both providers. Mean wall-clock per document
is 9.65~seconds for GPT-4o/Vision against 14.44~seconds for
GPT-4o/OCR ($-33\%$), and 12.16~seconds for Sonnet/Vision against
16.35~seconds for Sonnet/OCR ($-26\%$). Figure~\ref{fig:eval-timing}
splits the time. Tesseract accounts for $\sim$8.7~seconds of every
OCR-mediated run; LLM-call latency rises in vision mode (Sonnet/Vision
8.6~seconds vs.\ Sonnet/OCR 3.7~seconds; GPT-4o/Vision 6.2~seconds
vs.\ GPT-4o/OCR 2.1~seconds) because the model now does the page
reading internally, but removing Tesseract is the larger effect on
total time on this hardware.

\begin{figure}[ht]
    \centering
    \includegraphics[width=0.85\linewidth]{Chapters/figures/fig-7-6-timing-breakdown.png}
    \caption[Timing breakdown.]{Mean per-document timings split into
    OCR and LLM stages. Tesseract dominates the OCR-mediated
    pipeline on this hardware; the LLM-call cost rises in vision
    mode but does not close the gap.}
    \label{fig:eval-timing}
\end{figure}

The latency ranking is hardware-conditioned. Tesseract was run
single-process on commodity CPU and is embarrassingly parallel, so
the OCR-side latency would shrink substantially on a multi-core
deployment. The latency ranking above is therefore a single-laptop
ranking, not a production-scale verdict; the production-scale
prediction is that the LLM-call dominates wall-clock under either
modality once OCR is parallelised, and the cross-modality latency
gap closes.

\subsection{Operational complexity}
\label{sec:eval-rq2-complexity}

Vision is simpler to deploy than OCR, and the ranking is
unambiguous. The OCR-mediated pipeline depends on Poppler (system
binary), Tesseract (system binary plus per-language data files for
\texttt{eng}, \texttt{deu}, \texttt{fra} and \texttt{ita}) and the
EU/US-aware float parser of
\S\ref{sec:design-normalisation}. The vision-mediated pipeline
depends on the provider SDK alone. Adding a fifth language to the
OCR-mediated pipeline requires a new Tesseract data file on every
host; adding a fifth language to the vision-mediated pipeline
requires nothing.

The H2 verdict is therefore mixed: vision is more expensive but
faster and operationally simpler. The trade-off direction H2
predicted is right; the per-axis directions split along provider
lines and along the cost-vs-latency axis in ways the as-written
hypothesis did not anticipate.

% ------------------------------------------------------------
\section{RQ3 --- does the choice of LLM matter?}
\label{sec:eval-rq3}

\subsection{Modality-by-provider interaction}
\label{sec:eval-rq3-interaction}

H3 (the modality gap differs by provider) is supported. The
within-provider modality gap is $+2.10\,\mathrm{pp}$ for Sonnet
and $+1.02\,\mathrm{pp}$ for GPT-4o, a gap-of-gaps of
$1.08\,\mathrm{pp}$. The interaction has a mechanistic reading:
Sonnet has more material to lose in vision mode because it is the
better OCR-mode reader on this corpus, and Sonnet's vision-mode
disadvantage concentrates on the M1 (header anchoring) and M3
(locale formatting) mechanisms above. GPT-4o is less prone to
header anchoring on the M1 cases inspected, partly because its
visible behaviour is to omit a field rather than substitute a
salient nearby string when the answer is uncertain
(\S\ref{sec:eval-rq3-orthogonal}).

\subsection{Provider effects orthogonal to modality}
\label{sec:eval-rq3-orthogonal}

Two provider effects appear in both modalities and are therefore
attributable to the model rather than to the input mode.

The first is currency contextual inference. The four BC~Hydro bills
have \texttt{currency} ground truth \texttt{CAD}, but the string
``CAD'' does not appear in the OCR text or on the page image: the
bills display the bare ``\$'' sign. Sonnet returns \texttt{CAD} on
all four under both OCR and vision; GPT-4o returns \texttt{USD} on
all four under both modalities. The signal Sonnet uses (``BC
Hydro'', ``Vancouver'', Canadian postal code) is available to both
providers; GPT-4o reads it and does not act on it. This single
mechanism accounts for the entire 0.90 vs.\ 1.00 currency gap in
Table~\ref{tab:eval-headline} between GPT-4o and Sonnet.

The second is the hallucination-versus-omission profile. Across
both modalities, Sonnet emits 20 hallucinations (8 OCR, 12 vision)
and 20 omissions (9 OCR, 11 vision); GPT-4o emits 5 hallucinations
(1 OCR, 4 vision) and 32 omissions (17 OCR, 15 vision). Sonnet
hallucinates roughly four times as often as GPT-4o; GPT-4o omits
where Sonnet guesses. Figure~\ref{fig:eval-pareto} plots the four
conditions on the two axes and shows that the providers occupy
distinct points, not different points on a single curve. Downstream
systems that punish wrong values should prefer GPT-4o; ones that
punish missing values should prefer Sonnet.

\begin{figure}[ht]
    \centering
    \includegraphics[width=0.78\linewidth]{Chapters/figures/fig-7-9-hallucination-omission-pareto.png}
    \caption[Hallucination-versus-omission Pareto.]{Per-condition
    counts of hallucinations and omissions across the corpus.
    Sonnet sits in the high-hallucination low-omission region under
    both modalities; GPT-4o sits in the low-hallucination
    higher-omission region. The two providers occupy distinct
    operating points, not different points on a single curve.}
    \label{fig:eval-pareto}
\end{figure}

\subsection{Cross-provider agreement and ensemble ceiling}
\label{sec:eval-rq3-agreement}

The two providers agree on the same answer for $\sim$80.4\% of OCR-mode
field comparisons and $\sim$78.1\% of vision-mode comparisons, where
``agree'' means both predictions normalise to the same value
regardless of correctness. Disagreements concentrate on the
mechanisms above: BC~Hydro \texttt{currency}, German credit-bill
totals, the Affinity Water \texttt{total\_amount\_due}
disambiguation, and \texttt{provider\_name} on bills with strong
brand wordmarks. A field-level vote ensemble of Sonnet and GPT-4o
in OCR mode would clear $\sim$91.4\% accuracy on this corpus
(against the better of the two cells alone, 89.65\% for Sonnet/OCR);
in vision mode the corresponding upper bound is $\sim$89.6\% against
87.55\% for Sonnet/Vision. The provider-level personality difference
H3 predicts is therefore exploitable, but the ceiling is not large
enough to outweigh the cost of running both providers per document.

% ------------------------------------------------------------
\section{Where failures actually come from}
\label{sec:eval-failures}

\subsection{The diagnosis heuristic and what it misses}
\label{sec:eval-failures-heuristic}

The diagnosis heuristic specified in
\S\ref{sec:design-evalframework-diagnosis} attributes a non-correct
OCR-mediated field to OCR if the ground-truth value is not a
case-insensitive substring of the stored OCR text, and to the LLM
otherwise. The heuristic is reproducible and inspectable, and it is
the right starting point. It also has known failure modes that
matter for the analysis below.

The first is contextual inference. On the BC~Hydro currency cases
above, the ground-truth string ``CAD'' is not in the OCR text, so
the strict heuristic attributes the GPT-4o failure to OCR. The
actual cause is that GPT-4o did not infer the currency from the
issuer name and address; Sonnet did infer it from exactly the same
OCR text. Attributing this to OCR is wrong in the sense that no OCR
change would have fixed it. The second is ground-truth phrasing: on
\texttt{us\_electricity\_publicservicecompany\_001} the GT
\texttt{provider\_name} is the legal name ``Public Service Company
of Colorado'', and the OCR text contains the brand wordmark ``Xcel
Energy'' rather than the legal name. The strict heuristic attributes
this to OCR; the actual cause is annotation choosing the legal name
where the brand wordmark would have been an equally defensible
answer. The third is locale-formatted numerics: on a bill that
displays ``\texttt{16.250}'' the GT records ``\texttt{16250}''
(thousands separator removed) and the heuristic again says OCR even
though the text was readable.

\subsection{Two-bucket vs.\ four-bucket re-attribution}
\label{sec:eval-failures-fourbucket}

Replacing two buckets with four addresses the failure modes above.
The four buckets are: \emph{Pure OCR} (Tesseract did not capture
the value), \emph{Pure LLM-extraction} (the LLM mis-read text that
was present), \emph{LLM-inference} (context for inference was
present and the model did not act on it; the BC~Hydro currency
cases sit here), and \emph{GT-ambiguity} (the annotation chose one
of multiple defensible answers; the PSCo and Affinity Water cases
sit here). The re-classification rule is: known multiple-answer
cases move to GT-ambiguity, currency and \texttt{utility\_type}
strict-OCR-failures move to LLM-inference, the rest stay where the
strict heuristic put them. Figure~\ref{fig:eval-attribution} shows
the result.

\begin{figure}[ht]
    \centering
    \includegraphics[width=0.95\linewidth]{Chapters/figures/fig-7-10-failure-attribution.png}
    \caption[Failure attribution: two-bucket vs.\ four-bucket.]{Failure
    attribution for the two OCR-mediated runs. Panel (a): the
    strict two-bucket heuristic of
    \S\ref{sec:design-evalframework-diagnosis}. Panel (b): four
    buckets after applying the re-classification rule of
    \S\ref{sec:eval-failures-fourbucket} to
    \texttt{derived/diagnosis\_audit.csv}.}
    \label{fig:eval-attribution}
\end{figure}

Of 47 Sonnet/OCR non-correct fields, 26 (55\%) are Pure OCR, 17
(36\%) are Pure LLM-extraction, 1 is LLM-inference and 3 are
GT-ambiguity. Of 61 GPT-4o/OCR non-correct fields, 31 (51\%) are
Pure OCR, 23 (38\%) are Pure LLM-extraction, 4 are LLM-inference
and 3 are GT-ambiguity. The headline finding survives in spirit: a
little over half of OCR-mode failures are addressable by improving
Tesseract; the remainder are addressable by improving the LLM, the
prompt or the annotation. The four-bucket reading sharpens that
finding by separating the cases where the OCR text was the problem
from the cases where the OCR text was fine and the model did not
act on what it had.

\subsection{Modality-specific patterns}
\label{sec:eval-failures-modality}

Two failure types appear under one input mode and not the other.
The OCR-mediated path hallucinates redacted spans into plausible
characters: on
\texttt{it\_electricity\_irenmercato\_001} the ground-truth
\texttt{bill\_number} is \texttt{502305934151}; the OCR text
contains the redacted form ``\texttt{50230XXXXXXX}'' and the LLM
returns a guessed completion in OCR mode while correctly returning
\texttt{null} (or the unredacted prefix) in vision mode where the
redaction box is visible as a layout element. The vision-mediated
path inverts the failure mode: on bills with prominent brand
wordmarks the model anchors to the wordmark on
\texttt{provider\_name} (M1), and on bills with European-locale
thousands separators the model collapses the separator (M3).

A third pattern affects every condition equally: sign handling on
credit and refund bills. On \texttt{de\_electricity\_naturstromhandel\_001}
the ground-truth \texttt{total\_amount\_due} is $-40.14$~EUR
because the bill is a refund (``Guthaben''). All four conditions
return $40.14$, dropping the sign. This is not a modality
strength on either side; it is a construct issue between the prompt
and the ground-truth convention, and it is recorded as a
limitation in \S\ref{sec:eval-validity}.

\subsection{The two-page baseline as a deliberate design choice}
\label{sec:eval-failures-pagecap}

Both pipelines process only the first two pages of each bill. The
cap is set explicitly in \texttt{src/pipeline.py} (lines 137 and
149: \texttt{ocr\_images = images[:2]} and
\texttt{vision\_images = images[:2]}). It is a deliberate baseline,
not a confound. On every bill in this corpus, every required field
that the bill carries appears within the first two pages: the
header (provider, utility type, bill number, dates), the totals
panel and the consumption box are universally on page one or two,
and later pages carry per-meter detail, terms and conditions,
multilingual translations or marketing inserts that the schema does
not target. The cap therefore costs almost nothing in accuracy
under this corpus while cutting input tokens by a factor of
$\sim$3.9 (the mean page count is 3.9 with the longest bill at
nine pages) and removing $\sim$5--7 seconds of Tesseract time on
the OCR side per long bill. Figure~\ref{fig:eval-page-cap}
visualises which documents are seen in full and which are
truncated.

\begin{figure}[ht]
    \centering
    \includegraphics[width=0.78\linewidth]{Chapters/figures/fig-7-8-page-cap-coverage.png}
    \caption[Two-page cap coverage.]{Per-document coverage under
    the two-page cap, applied to both pipelines. Dark segments are
    the pages either pipeline reads; light segments are the pages
    neither pipeline reads.}
    \label{fig:eval-page-cap}
\end{figure}

The cap is a property of the input distribution, not of the
algorithm. On a corpus that exercises later-page targets, changing
the slice from \texttt{[:2]} to a higher value, or to a function
that picks pages by a header-detection heuristic, would extend
coverage at proportional cost. The cost-vs-coverage curve for that
extension is the most direct empirical follow-up the corpus
suggests, and is left for Chapter~\ref{chap:conclusion}. The
relevant point here is that the cap is the reason cost and latency
are flat in document length on this corpus, and the reason the
modality-vs-modality comparison in \S\ref{sec:eval-rq1} is a
like-for-like comparison rather than a coverage-difference artefact.

% ------------------------------------------------------------
\section{Counterfactuals against the alternative designs}
\label{sec:eval-alternatives}

This section returns to the alternatives recorded in
Figure~\ref{fig:design-decisions} and tests them against the run
evidence rather than against intuition; the analysis closes the
A++ requirement that an evaluation chapter compare results expected
from theory for alternative designs with those obtained in
practice (Chapter~\ref{chap:spec}~\S\ref{sec:spec-evaluation}).

\subsection{Per-issuer templates}
\label{sec:eval-alternatives-templates}

The data ratifies the rejection of \S\ref{sec:design-decisions}.
Thirty-three unique issuers appear in the corpus and the field-level
error patterns are not positional: \texttt{provider\_name} fails on
brand-vs-legal selection (M1), \texttt{consumption\_amount} fails on
locale-aware separator handling (M3), and the
\texttt{total\_amount\_due} sign loss on credit bills is a semantic
question, not a coordinate-extraction question. None of those failure
modes is template-fixable; each one requires either an
issuer-specific rule or a flexible reading of the bill, and the
LLM-mediated approach gives the latter for free.

\subsection{Fine-tuned encoder--decoder (LayoutLMv3, Donut)}
\label{sec:eval-alternatives-finetune}

Forty documents is well below the floor at which Donut and
LayoutLMv3 are reported usable~\cite{kim2022donut,huang2022layoutlmv3};
provider concentration would inject training-distribution bias the
holdout cannot detect (BC~Hydro at four documents would be
disproportionately easy on any random split); and per-language
fine-tuning would be needed for the Italian and German subsets where
mechanisms M3 and the credit-bill sign issue occur. The LLM-mediated
approach pays a per-call cost in exchange for zero training data,
and at \$0.008--0.015 per document the per-call cost is below the
cost of one annotation hour at any reasonable rate.

\subsection{Provider ensembling}
\label{sec:eval-alternatives-ensemble}

The ensembling alternative is testable directly against the run
evidence (\S\ref{sec:eval-rq3-agreement}). A field-level vote of
Sonnet/OCR and GPT-4o/OCR clears $\sim$91.4\% on this corpus,
against 89.65\% and 86.35\% for the two cells alone. Most of the
gain is on the disagreement set (BC~Hydro \texttt{currency},
PSCo \texttt{provider\_name}, the German credit-bill totals where
Sonnet recovers the sign for one of the four cases). The ceiling
is real but moderate; pairing it with the cost ranking in
\S\ref{sec:eval-rq2-cost} (running two providers per document
roughly doubles the cost) puts the operating decision squarely in
the downstream-cost-tolerance domain.

\subsection{Single-modality system}
\label{sec:eval-alternatives-singlemod}

A vision-only or OCR-only system would have to cope with the
within-document modality flipping documented in
\S\ref{sec:eval-rq1-flips}. Twelve documents per provider have at
least one field where one modality is correct and the other wrong,
and several of those documents the directions oppose. A perfect
modality router would beat the better of the two single-modality
cells by roughly the same amount as a provider ensemble. The case
for keeping both pipelines is precisely this: not because either
wins outright, but because the conditions under which each is
preferable are field-, locale- and issuer-specific.

\subsection{Vision (or OCR) with all pages}
\label{sec:eval-alternatives-allpages}

The two-page cap is the design choice that makes both modalities
flat in document length on this corpus
(\S\ref{sec:eval-failures-pagecap}). Removing the cap on a corpus
where every required field already lives on the first two pages
would buy nothing in accuracy and cost roughly $3.9\times$ more in
input tokens (and proportionally more in OCR latency) on average.
Removing the cap on a different corpus where required fields appear
on later pages would buy accuracy at proportional cost, and the
cost-vs-coverage curve as a function of page count is the
empirical question Chapter~\ref{chap:conclusion} picks up. Either
way, the cap is the right choice for the corpus this study
evaluates.

% ------------------------------------------------------------
\section{What these numbers cannot say}
\label{sec:eval-validity}

\subsection{Internal validity}
\label{sec:eval-validity-internal}

Two internal threats matter. The diagnosis heuristic mis-attributes
the cases listed in \S\ref{sec:eval-failures-heuristic}; the
four-bucket re-classification of
\S\ref{sec:eval-failures-fourbucket} is the corrective. Sign
handling on credit-bill totals fails uniformly across all four
conditions because the prompt does not encode the
``\texttt{total\_amount\_due} is signed (negative on credit
bills)'' convention that the ground truth uses; this is a prompt
construct issue rather than a model failure, and it bounds the
\texttt{total\_amount\_due} accuracy ceiling on this corpus to
roughly 95\% with the current prompt.

\subsection{Construct validity}
\label{sec:eval-validity-construct}

Three documented cases of multiple-defensible-answer GT. On
\texttt{gb\_water\_affinitywater\_001} the GT \texttt{total\_amount\_due}
is \pounds46.00 (current charges); the bill also displays
\pounds235.14 (outstanding balance). On
\texttt{us\_gas\_peoplesgas\_001} the GT \texttt{account\_number}
is recorded with a particular whitespace pattern; the bill displays
the same digits with a different whitespace pattern. On
\texttt{us\_electricity\_publicservicecompany\_001} the GT
\texttt{provider\_name} is the legal entity ``Public Service Company
of Colorado''; the bill prominently displays ``Xcel Energy'', the
parent brand. The annotator picked one of two valid answers in
each case. Under the alternative defensible conventions the headline
numbers move by a few tenths of a percentage point and the modality
and provider rankings do not change.

\subsection{External validity and acknowledged limitations}
\label{sec:eval-validity-external}

Six bounds on external validity, each with its cause and the
specific finding it bounds.

\textbf{Corpus size.} $n=40$ active documents. Per-language and
per-utility-type slices are descriptive at this size; differences
of a few percentage points should not be over-read. The cause is
the manual collection constraint of
\S\ref{sec:eval-dataset-provenance}; expansion to $n\geq 200$
would lift the slice analyses to formally testable form.

\textbf{Language and script imbalance.} 20~English, 12~German,
7~Italian, 1~French; Latin script only. Non-Latin scripts (Arabic,
Cyrillic, CJK) are out of scope. The per-language results in
\S\ref{sec:eval-rq1-stratified} should not be extrapolated to
language families not exercised here.

\textbf{Single-annotator labelling.} Inter-annotator agreement
cannot be computed from a single labelling pass. The construct
threats above are the visible price; the multiple-defensible-answer
cases were caught by the author during error analysis but a second
pass would catch more.

\textbf{Single OCR engine.} The OCR-mediated pipeline uses
Tesseract only. Alternative OCR engines (Google Document AI,
PaddleOCR, AWS Textract) carry different error profiles, and the
``$\sim$50\% of OCR-mode failures are pure OCR'' finding is
therefore engine-specific. A drop-in replacement of the
\texttt{OCREngine} class (\S\ref{sec:design-pipeline}) is the
intended extension surface; running it remains future work.

\textbf{Two LLM providers, one model each.} Anthropic Sonnet 4.5
and OpenAI GPT-4o are the only two models tested. The
hallucination-vs-omission profile in \S\ref{sec:eval-rq3-orthogonal}
is therefore a two-point sample, not a population claim. Adding
Google Gemini, Mistral Large or open-weight models such as
Llama~3.2 Vision would broaden the provider axis and is the second
intended extension surface.

\textbf{Born-digital corpus only.} No scanned bills are included.
The OCR error ceiling identified in Lopresti's
analysis~\cite{lopresti2009ocrerrors} is therefore not exercised
at its hardest, and the OCR side of the modality comparison is
shown at its strongest. A follow-up pass with scanned reproductions
would isolate the scan-quality effect against the findings here.

% ------------------------------------------------------------
\section{Synthesis}
\label{sec:eval-synthesis}

Five findings carry forward to Chapter~\ref{chap:conclusion}.

\textbf{H1 is rejected on this corpus.} Vision-mediated extraction
did not match OCR-mediated extraction in aggregate, and the gap is
provider- and field-specific rather than uniform.

\textbf{H2 partially holds.} Vision is more expensive (consistent
with H2) and operationally simpler than OCR; vision is faster
end-to-end (against H2), and the latency reading is conditional on
single-process Tesseract.

\textbf{H3 is supported.} The within-provider modality gap differs
by provider, the providers occupy distinct points in
hallucination-versus-omission space, and the modality choice
interacts with provider choice rather than being orthogonal to it.

\textbf{Modality is a per-field decision.} Within-document modality
flipping (Figure~\ref{fig:eval-flips}) implies that no
single-modality system is uniformly best on this corpus, and the
empirical case for keeping both pipelines is the field-level
asymmetry, not aggregate accuracy.

\textbf{Half of OCR-mode failures are not Tesseract's fault.} The
four-bucket re-attribution shows that 45--49\% of OCR-mediated
failures are addressable by improving the LLM, the prompt or the
annotation rather than the preprocessing.

These five statements are the substrate Chapter~\ref{chap:conclusion}
turns into transferable claims and predictions for new domains.
